\subsection{扩张模与扭积的对偶性}

设 A 为环，P 与 J 为 A-模。对每个 A-模复形 C，我们在 A, X, 第 99 页，命题 12 中构造了复形的一个典范同构

\[
\mu : \operatorname{Homgr} _ {\mathrm{A}} (\mathrm{C} \otimes_ {\mathrm{A}} \mathrm{P}, \mathrm{J}) \longrightarrow \operatorname{Homgr} _ {\mathrm{A}} (\mathrm{C}, \operatorname{Hom} _ {\mathrm{A}} (\mathrm{P}, \mathrm{J})) .
\]

设 M 为\(\Lambda\) 模，(C, p) 为 M 的一个投射分解。考虑同态序列

\[
\begin{array}{r l} \operatorname{Ext} _ {\mathrm{A}} (\mathrm{M}, \operatorname{Hom} _ {\mathrm{A}} (\mathrm{P}, \mathrm{J})) & \xrightarrow {\varphi^ {- i}} \mathrm{H} (\operatorname{Homgr} _ {\mathrm{A}} (\mathrm{C}, \operatorname{Hom} _ {\mathrm{A}} (\mathrm{P}, \mathrm{J}))) \xrightarrow {\mathrm{H} (\mu) ^ {- 1}} \mathrm{H} (\operatorname{Homgr} _ {\mathrm{A}} (\mathrm{C} \otimes_ {\mathrm{A}} \mathrm{P}, \mathrm{J})) \\ & \xrightarrow {u} \operatorname{Homgr} _ {\mathrm{A}} (\mathrm{H} (\mathrm{C} \otimes_ {\mathrm{A}} \mathrm{P}), \mathrm{J}) \xrightarrow {v} \operatorname{Homgr} _ {\mathrm{A}} (\operatorname{Tor} ^ {\mathrm{A}} (\mathrm{M}, \mathrm{P}), \mathrm{J}), \end{array}
\]

其中 \(\varphi\) 是典范同构 \(\varphi(\mathrm{C},\mathrm{Hom}_{\mathrm{A}}(\mathrm{P},\mathrm{J}))\)（A, X, 第 100 页，定理 1），\(u\) 是典范同态 \(\lambda(\mathrm{C}\otimes_{\Lambda}\mathrm{P},\mathrm{J})\)（A, X, 第 82 页），而 \(v\) 由典范同构 \(\psi(\mathrm{C},\mathrm{P}):\operatorname{Tor}^{\Lambda}(\mathrm{M},\mathrm{P})\longrightarrow\mathrm{H}(\mathrm{C}\otimes_{\mathrm{A}}\mathrm{P})\) 导出。

设 \((\mathbf{C}', p')\) 为 M 的另一个投射分解。由 A, X, 第 49 页，命题 3 的推论，存在复形的一个同伦 \(\alpha: \mathbf{C}' \to \mathbf{C}\)，使 \(p \circ \alpha = p'\)。由 A, X, 第 103 页，命题 2 可知，对每个 A-模 R，有 \(\mathrm{H}(\alpha \otimes 1_{\mathrm{P}}) \circ \psi(\mathrm{C}', \mathrm{P}) = \psi(\mathrm{C}, \mathrm{P})\) 与 \(\varphi(\mathrm{C}', \mathrm{R}) \circ \mathrm{H}(\mathrm{Homgr}(\alpha, 1_{\mathrm{R}})) = \varphi(\mathrm{C}, \mathrm{R})\)。由此推出，0 次分次同态

\[
\theta (\mathrm{M}, \mathrm{P}): \operatorname{Ext} _ {\mathrm{A}} (\mathrm{M}, \operatorname{Hom} _ {\mathrm{A}} (\mathrm{P}, \mathrm{J})) \longrightarrow \operatorname{Homgr} _ {\mathrm{A}} \left(\operatorname{Tor} ^ {\mathrm{A}} (\mathrm{M}, \mathrm{P}), \mathrm{J}\right)
\]

（即上述同态序列的合成）与 M 的投射分解 (C, p) 的选取无关。按构造，它是 End \(_{A}\) (J)-线性的。

同态 \(\theta(\mathrm{M},\mathrm{P})\) 的定义可明确表述如下。设 \(p\) 为整数，\(\mathfrak{v}\) 为 \(\operatorname{Ext}_{\Lambda}^{p}(\mathrm{M},\operatorname{Hom}_{\mathrm{A}}(\mathrm{P},\mathrm{J}))\) 的一个元素，\(\tau\) 为 \(\operatorname{Tor}_{p}^{\mathrm{A}}(\mathrm{M},\mathrm{P})\) 的一个元素。利用同构 \(\varphi(\mathrm{C},\operatorname{Hom}_{\mathrm{A}}(\mathrm{P},\mathrm{J}))\)，\(\mathfrak{v}\) 由线性映射 \(u:\mathrm{C}_{p}\to\operatorname{Hom}_{\mathrm{A}}(\mathrm{P},\mathrm{J})\) 代表，且满足 \(u\circ d_{\mathrm{C}}=0\)；类似地，利用 \(\psi(\mathrm{C},\mathrm{P})\)，\(\tau\) 由元素 \(\sum c_{\mu}\otimes p_{\mu}\)（属于 \(\mathrm{C}_{p}\otimes\mathrm{P}\)，满足 \(\sum d_{\mathrm{C}}(c_{\mu})\otimes p_{\mu}=0\)）代表。于是有 \(\theta(\mathrm{M},\mathrm{P})(\mathfrak{v})(\tau)=\sum u(c_{\mu})(p_{\mu})\)。

另一方面，设 \(\nu : C \otimes_A \operatorname{Hom}_A(P, J) \longrightarrow \operatorname{Homgr}_A(\operatorname{Homgr}_A(C, P), J)\) 为把元素 \(c \otimes h\)（其中 \(c \in C_p\)，\(h \in \operatorname{Hom}_A(P, J)\)）映到同态 \(u \mapsto (-1)^p h(u(c))\) 的同态。它是 0 次分次的；若每个模 \(C_p\) 都是有限型自由模，则它是双射。不难验证它是复形的一个态射。

考虑同态序列

\[
\begin{array}{r l} \operatorname{Tor} ^ {\Lambda} (M, \operatorname{Hom} _ {\Lambda} (P, J)) & \xrightarrow {\psi} H (C \otimes_ {A} \operatorname{Hom} _ {\Lambda} (P, J)) \xrightarrow {H (\nu)} H (\operatorname{Homgr} _ {A} (\operatorname{Homgr} _ {A} (C, P), J)) \\ & \xrightarrow {w} \operatorname{Homgr} _ {\Lambda} (H (\operatorname{Homgr} _ {A} (C, P)), J) \xrightarrow {t} \operatorname{Homgr} _ {A} (\operatorname{Ext} _ {A} (M, P), J) \end{array}
\]

其中 \(\psi\) 是典范同构 \(\psi(\mathrm{C},\mathrm{Hom}_{\mathrm{A}}(\mathrm{P},\mathrm{J}))\)，\(w\) 是典范同态 \(\lambda(\mathrm{Homgr}_{\mathrm{A}}(\mathrm{C},\mathrm{P}),\mathrm{J})\)（A, X, 第 82 页），而 \(t\) 由典范同构 \(\varphi(\mathrm{C},\mathrm{P})\) 导出。与上面一样可以看出，合成同态

\[
\rho (\mathrm{M}, \mathrm{P}): \operatorname{Tor} ^ {\mathrm{A}} (\mathrm{M}, \operatorname{Hom} _ {\mathrm{A}} (\mathrm{P}, \mathrm{J})) \longrightarrow \operatorname{Homgr} _ {\mathrm{A}} \left(\operatorname{Ext} _ {\mathrm{A}} (\mathrm{M}, \mathrm{P}), \mathrm{J}\right)
\]

与分解 C 的选取无关；它是 \(\operatorname{End}_{\mathrm{A}}(\mathrm{J})\)-线性的。设 p 为整数，\(\xi \in \operatorname{Tor}_{p}^{\mathrm{A}}(\mathrm{M}, \operatorname{Hom}_{\mathrm{A}}(\mathrm{P}, \mathrm{J}))\)，\(\lambda \in \operatorname{Ext}_{\mathrm{A}}^{p}(\mathrm{M}, \mathrm{P}), \mathrm{J}\)；若 \(\xi\) 借助 \(\psi(\mathrm{C}, \operatorname{Hom}_{\Lambda}(\mathrm{P}, \mathrm{J}))\) 由元素 \(\sum c_{\mu} \otimes u_{\mu}\)（属于 \(\mathrm{C} \otimes \operatorname{Hom}_{\Lambda}(\mathrm{P}, \mathrm{J})\)，满足 \(\sum d_{\mathrm{C}}(c_{\mu}) \otimes u_{\mu} = 0\)）代表，而 \(\lambda\) 借助 \(\varphi(\mathrm{C}, \mathrm{P})\) 由同态 \(\ell : C_{p} \to P\)（满足 \(\ell \circ d_{C} = 0\)）代表，则有 \(\rho(\mathrm{M}, \mathrm{P})(\xi)(\lambda) = (-1)^{p} \sum u_{\mu} (\ell(c_{\mu}))\)。

命题 6.—— 设 A-模 J 是内射的；对每个 A-模 N，令 D(N) = \(\mathrm{Hom}_{\Lambda}(N,J)\)。

\begin{enumerate}
\def\labelenumi{\alph{enumi})}
\item
  同态 \(\theta^{i}(\mathrm{M},\mathrm{P}):\operatorname{Ext}_{\mathrm{A}}^{i}(\mathrm{M},\mathrm{D}(\mathrm{P}))\longrightarrow \mathrm{D}\big(\operatorname {Tor}_{i}^{\mathrm{A}}(\mathrm{M},\mathrm{P})\big)\) 是双射。
\item
  若环 A 是 Noether 环且 A-模 M 是有限生成的，则同态 \(\rho_{i}(\mathrm{M},\mathrm{P}):\operatorname{Tor}_{i}^{\mathrm{A}}(\mathrm{M},\mathrm{D}(\mathrm{P}))\longrightarrow \mathrm{D}\bigl (\operatorname {Ext}_{\mathrm{A}}^{i}(\mathrm{M},\mathrm{P})\bigr)\) 是双射。
\item
  按构造，只要 \(\theta(\mathrm{M},\mathrm{P})\) 是双射，同态 \(\lambda(\mathrm{C}\otimes_{\mathrm{A}}\mathrm{P},\mathrm{J})\) 便是双射的；当 J 内射时即属此种情形（A, X, 第 85 页，推论 2）。
\item
  选取分解 C，使每个模 \(C_{p}\) 都是有限型自由模（A, X, 第 53 页，命题 6）。于是同态 v 是双射，\(\lambda(\text{Homgr}_{\text{A}}(\text{C}, \text{P}), \text{J})\) 也是双射，因为 J 是内射的；因此 \(\rho(\text{M}, \text{P})\) 是双射。
\end{enumerate}

注.—— 1) 对 A-模的每个同态 \(f: \mathrm{N} \to \mathrm{N}'\)，用 \(\mathrm{D}(f): \mathrm{D}(\mathrm{N}') \to \mathrm{D}(\mathrm{N})\) 表示同态 \(\mathrm{Hom}(f, \mathbf{1}_{\mathrm{J}})\)。设 \(u: \mathrm{M} \to \mathrm{M}'\) 与 \(v: \mathrm{P} \to \mathrm{P}'\) 为 A-模的同态。选取 M 的投射分解 (C, p) 与 M' 的投射分解 (C', p')，以及复形的一个态射 \(\tilde{u}: \mathrm{C} \to \mathrm{C}'\)，使 \(p' \circ \tilde{u} = u \circ p\)（A, X, 第 49 页，命题 3）。图表

\[
\begin{array}{c c c} \operatorname{Homgr} _ {\mathrm{A}} (\mathrm{C} ^ {\prime} \otimes_ {\mathrm{A}} \mathrm{P} ^ {\prime}, \mathrm{J}) & \xrightarrow {\mu^ {\prime}} & \operatorname{Homgr} _ {\mathrm{A}} (\mathrm{C} ^ {\prime}, \operatorname{Hom} _ {\mathrm{A}} (\mathrm{P} ^ {\prime}, \mathrm{J})) \\ \operatorname{Hom} (\tilde {u} \otimes v, 1 _ {\mathrm{J}}) \Bigg \downarrow & & \Bigg \downarrow \operatorname{Hom} (\tilde {u}, \operatorname{Hom} (v, 1)) \\ \operatorname{Homgr} _ {\mathrm{A}} (\mathrm{C} \otimes_ {\mathrm{A}} \mathrm{P}, \mathrm{J}) & \xrightarrow {\mu} & \operatorname{Homgr} _ {\mathrm{A}} (\mathrm{C}, \operatorname{Hom} _ {\mathrm{A}} (\mathrm{P}, \mathrm{J})) \end{array}
\]

其中 \(\mu\) 与 \(\mu'\) 为典范同态，该图是交换的；于是由 A, X, 第 103 页，命题 2 推出一个交换图

\[
\begin{array}{c c c} \operatorname{Ext} _ {\Lambda} ^ {i} (\mathrm{M} ^ {\prime}, \mathrm{D} (\mathrm{P} ^ {\prime})) & \xrightarrow {\theta^ {i} (\mathrm{M} ^ {\prime} , \mathrm{P} ^ {\prime})} & \mathrm{D} (\operatorname{Tor} _ {i} ^ {\mathrm{A}} (\mathrm{M} ^ {\prime}, \mathrm{P} ^ {\prime})) \\ \operatorname{Ext} ^ {i} (u, \mathrm{D} (v)) \Bigg \downarrow & & \Bigg \downarrow \mathrm{D} (\operatorname{Tor} _ {i} (u, v)) \\ \operatorname{Ext} _ {\Lambda} ^ {i} (\mathrm{M}, \mathrm{D} (\mathrm{P})) & \xrightarrow {\theta^ {i} (\mathrm{M} , \mathrm{P})} & \mathrm{D} (\operatorname{Tor} _ {i} ^ {\Lambda} (\mathrm{M}, \mathrm{P})) \end{array}
\]

设 \(w: \mathrm{P}'' \to \mathrm{P}\) 为 A-模的一个同态；以类似的方式得到一个交换图

\[
\begin{array}{c c c} \operatorname{Tor} _ {i} ^ {\mathrm{A}} (\mathrm{M}, \mathrm{D} (\mathrm{P})) & \xrightarrow {\rho_ {i} (\mathrm{M} , \mathrm{P})} & \mathrm{D} (\operatorname{Ext} _ {\mathrm{A}} ^ {i} (\mathrm{M}, \mathrm{P})) \\ \operatorname{Tor} _ {i} (u, \mathrm{D} (w)) \Bigg \downarrow & & \Bigg \downarrow \mathrm{D} (\operatorname{Ext} ^ {i} (u, w)) \\ \operatorname{Tor} _ {i} ^ {\mathrm{A}} (\mathrm{M} ^ {\prime}, \mathrm{D} (\mathrm{P} ^ {\prime \prime})) & \xrightarrow {\rho_ {i} (\mathrm{M} ^ {\prime} , \mathrm{P} ^ {\prime \prime})} & \mathrm{D} (\operatorname{Ext} _ {\mathrm{A}} ^ {i} (\mathrm{M} ^ {\prime}, \mathrm{P} ^ {\prime \prime})) \end{array} .
\]

\begin{enumerate}
\def\labelenumi{\arabic{enumi})}
\setcounter{enumi}{1}
\tightlist
\item
  设
\end{enumerate}

\[
0 \to \mathrm{M} ^ {\prime} \xrightarrow {j} \mathrm{M} \xrightarrow {q} \mathrm{M} ^ {\prime \prime} \to 0\tag{ℰ}
\]

为 A-模的一个正合列。在典范自由分解上诱导的同态 \(\mathrm{L}(q):\mathrm{L}(\mathrm{M})\to\mathrm{L}(\mathrm{M}^{\prime\prime})\) 是满射，而复形 \(\operatorname{Ker}\mathrm{L}(q)\) 给出 \(M^{\prime}\) 的一个投射分解。把 A, X, 第 104 页，命题 3 应用于正合列 \(0\to\operatorname{Ker}\mathrm{L}(q)\to\mathrm{L}(\mathrm{M})\to\mathrm{L}(\mathrm{M}^{\prime\prime})\to0\)，便得到交换图

\[
\begin{array}{c c c} \operatorname{Ext} _ {\mathrm{A}} ^ {i} (\mathrm{M} ^ {\prime}, \mathrm{D} (\mathrm{P})) & \xrightarrow {\theta^ {i} (\mathrm{M} ^ {\prime} , \mathrm{P})} & \mathrm{D} (\operatorname{Tor} _ {i} ^ {\mathrm{A}} (\mathrm{M} ^ {\prime}, \mathrm{P})) \\ \delta^ {i} (\mathcal {E}, \mathrm{D} (\mathrm{P})) \Bigg \downarrow & & \Bigg \downarrow (- 1) ^ {i + 1} \mathrm{D} (\partial_ {i + 1} (\mathcal {E}, \mathrm{P})) \\ \operatorname{Ext} _ {\mathrm{A}} ^ {i + 1} (\mathrm{M} ^ {\prime \prime}, \mathrm{D} (\mathrm{P})) & \xrightarrow {\theta^ {i + 1} (\mathrm{M} ^ {\prime \prime} , \mathrm{P})} & \mathrm{D} (\operatorname{Tor} _ {i + 1} ^ {\mathrm{A}} (\mathrm{M} ^ {\prime \prime}, \mathrm{P})) \end{array}
\]

\[
\begin{array}{c c c} \operatorname{Tor} _ {i + 1} ^ {\mathrm{A}} (\mathrm{M} ^ {\prime \prime}, \mathrm{D} (\mathrm{P})) & \xrightarrow {\rho_ {i + 1} (\mathrm{M} ^ {\prime \prime} , \mathrm{P})} & \mathrm{D} (\operatorname{Ext} _ {\mathrm{A}} ^ {i + 1} (\mathrm{M} ^ {\prime \prime}, \mathrm{P})) \\ \partial_ {i + 1} (\mathcal {E}, \mathrm{D} (\mathrm{P})) \Bigg \downarrow & & \Bigg \downarrow (- 1) ^ {i + 1} \mathrm{D} (\delta^ {i} (\mathcal {E}, \mathrm{P})) \\ \operatorname{Tor} _ {i} ^ {\mathrm{A}} (\mathrm{M} ^ {\prime}, \mathrm{D} (\mathrm{P})) & \xrightarrow {\rho_ {i} (\mathrm{M} ^ {\prime} , \mathrm{P})} & \mathrm{D} (\operatorname{Ext} _ {\mathrm{A}} ^ {i} (\mathrm{M} ^ {\prime}, \mathrm{P})) \end{array}
\]

设

\[
0 \to \mathrm{P} ^ {\prime} \to \mathrm{P} \to \mathrm{P} ^ {\prime \prime} \to 0\tag{F}
\]

为 A-模的一个正合列；由于 A-模 J 是内射的，由此导出一个正合列

\[
0 \to \mathrm{D} (\mathrm{P} ^ {\prime \prime}) \to \mathrm{D} (\mathrm{P}) \to \mathrm{D} (\mathrm{P} ^ {\prime}) \to 0.\tag{\(\left(\mathcal{D}(\mathcal{F})\right)\)}
\]

把 A, X, 第 104 页，命题 3 以及第 106 页，命题 4 应用于正合列 (F) 与 (D(F))，便以类似的方式得到交换图。

\[
\begin{array}{c c c} \operatorname{Ext} _ {\Lambda} ^ {i} (\mathrm{M}, \mathrm{D} (\mathrm{P} ^ {\prime})) & \xrightarrow {\theta^ {i} (\mathrm{M} , \mathrm{P} ^ {\prime})} & \mathrm{D} (\operatorname{Tor} _ {i} ^ {\mathrm{A}} (\mathrm{M}, \mathrm{P} ^ {\prime})) \\ \delta^ {i} (\mathrm{M}, \mathrm{D} (\mathcal {F})) \Bigg \downarrow & & \Bigg \downarrow (- 1) ^ {i + 1} \mathrm{D} (\partial_ {i + 1} (\mathrm{M}, \mathcal {S})) \\ \operatorname{Ext} _ {\mathrm{A}} ^ {i + 1} (\mathrm{M}, \mathrm{D} (\mathrm{P} ^ {\prime \prime})) & \xrightarrow {\theta^ {i + 1} (\mathrm{M} , \mathrm{P} ^ {\prime \prime})} & \mathrm{D} (\operatorname{Tor} _ {i + 1} ^ {\mathrm{A}} (\mathrm{M}, \mathrm{P} ^ {\prime \prime})) \end{array}
\]

\[
\begin{array}{c c c} \operatorname{Tor} _ {i + 1} ^ {\mathrm{A}} (\mathrm{M}, \mathrm{D} (\mathrm{P} ^ {\prime})) & \xrightarrow {\rho_ {i + 1} (\mathrm{M} , \mathrm{P} ^ {\prime})} & D (\operatorname{Ext} _ {\Lambda} ^ {i + 1} (\mathrm{M}, \mathrm{P} ^ {\prime})) \\ \partial_ {i + 1} (\mathrm{M}, \mathrm{D} (\mathcal {F})) \Bigg \downarrow & & \Bigg \downarrow (\sim 1) ^ {i} \mathrm{D} (\delta^ {i} (\mathrm{M}, \mathcal {F})) \\ \operatorname{Tor} _ {i} ^ {\mathrm{A}} (\mathrm{M}, \mathrm{P} ^ {\prime \prime}) & \xrightarrow {\rho_ {i} (\mathrm{M} , \mathrm{P} ^ {\prime \prime})} & D (\operatorname{Ext} _ {\mathrm{A}} ^ {i} (\mathrm{M}, \mathrm{P} ^ {\prime \prime})) \end{array}
\]

\subsection{对偶化模}

\subsection{对偶化模}

定义 1.—— 设 A 为 Noether 环。称 Λ-模 Ω 为对偶化的，如果它是有限生成的，并且对 Λ 的每个极大理想 m，A/m-向量空间 Ext \(_{A}^{i}\) (A/m, Ω) 在 i ≠ ht(m) 时为零，而在 i = ht(m) 时维数为 1。

对 A 的每个极大理想 m 与每个整数 i，A/m-向量空间 \(\operatorname{Ext}_{\mathrm{A}}^{i}(\mathrm{A}/\mathfrak{m},\Omega)\) 典范地同构于 \(\operatorname{Ext}_{\mathrm{A}_{\mathfrak{m}}}^{i}(\mathrm{A}/\mathfrak{m},\Omega_{\mathfrak{m}})\)。因此，为使有限生成 A-模 \(\Omega\) 是对偶化的，必须且只须 \(A_{m}\)-模 \(\Omega_{m}\) 对 A 的每个极大理想 m 都是对偶化的。

例.—— 1) 若环 A 是局部 Artin 环，则对偶化 A-模就是有限生成的内射 A-模 \(\Omega\)，它使得 \(\mathrm{Hom}_{\mathrm{A}}(\kappa_{\mathrm{A}},\Omega)\)（记之为 \( E \)）的维数等于 1（\(\S\) 3, 第 3 节，命题 6），亦即 Matlis A-模（\(\S\) 8, 第 3 节）。

\begin{enumerate}
\def\labelenumi{\arabic{enumi})}
\setcounter{enumi}{1}
\tightlist
\item
  为使 Noether 环 A 是 Gorenstein 环，必须且只须 A-模 A 是对偶化的（§ 3, 第 7 节，命题 11）。特别地，当 A 正则时，A-模 A 是对偶化的。
\end{enumerate}

注.—— 1) 设 A 为 Noether 局部环，\(\Omega\) 为有限生成 A-模。剩余域 \(\kappa_{\widehat{A}}\) 等同于 \(\kappa_{A}\)，而 \(\widehat{A}\)-模 \(\widehat{\Omega}\) 等同于 \(\widehat{A} \otimes_{A} \Omega\)（III, § 3, 第 4 节，定理 3）。于是由 A, X, 第 111 页，命题 10 可知，\(\kappa_{A}\)-向量空间 \(\mathrm{Ext}_{\mathbf{A}}^{i}(\kappa_{\mathbf{A}}, \Omega)\) 典范地同构于 \(\mathrm{Ext}_{\widehat{A}}^{i}(\kappa_{\widehat{A}}, \widehat{\Omega})\)。因此，为使 \(A\)-模 \(\Omega\) 是对偶化的，必须且只须 \(\widehat{A}\)-模 \(\widehat{\Omega}\) 是对偶化的。

\begin{enumerate}
\def\labelenumi{\arabic{enumi})}
\setcounter{enumi}{1}
\tightlist
\item
  设 \(\Omega\) 为对偶化 A-模；对每个秩 1 的投射 A-模 L，A-模 \(\Omega\otimes_{\mathrm{A}}\mathrm{L}\) 是对偶化的（A, X, 第 108 页，命题 7, b)）。我们将在下面（\(n^{\circ}\) 4, 命题 6）看到，每个对偶化 A-模都同构于这种形状的模。
\end{enumerate}

命题 1.—— 设 A 为 Noether 环，Ω 为对偶化 A-模。

\begin{enumerate}
\def\labelenumi{\alph{enumi})}
\tightlist
\item
  则 A 是 Macaulay 环，且 A-模 Ω 是 Macaulay 模。
\end{enumerate}

\[
\text { b) } O n a \mathrm{di} _ {\mathrm{A}} (\Omega) = \dim (\Omega) = \dim (\mathrm{A}).
\]

先设环 A 是局部环，并用 \(d\) 表示它的维数。§ 3, 第 3 节的命题 6 蕴含 \(\mathrm{di}_{\Lambda}(\Omega) = d\)，从而由 § 3, 第 6 节的命题 9 得 \(\mathrm{prof}(\mathrm{A}) = d\)，故 A 是 Macaulay 环。此外，按深度的定义，有 \(\mathrm{prof}(\Omega) = d\)；由于有 \(\mathrm{prof}(\Omega) \leqslant \dim(\Omega) \leqslant d\)，由此即得此情形下的命题。

在一般情形，\(A_{m}\)-模 \(\Omega_{m}\) 对 A 的每个极大理想 m 都是对偶化的，从而 \(A_{m}\) 是 Macaulay 环，而 \(\Omega_{m}\) 是 \(\Lambda_{m}\)-模 a)。此外，对每个极大理想 m，有 \(\mathrm{di}_{\mathrm{A}_{\mathfrak{m}}(\Omega_{\mathfrak{m}})} = \dim(\Omega_{\mathfrak{m}}) = \dim(\Lambda_{\mathfrak{m}})\)，于是对 b) 取上确界即得（§ 3, 第 2 节，命题 3）。

命题 2.—— 设 A 为 Noether 环，Ω 为对偶化 A-模。对 A 的每个素理想 p，A\(_{p}\)-模 Ω\(_{p}\) 是对偶化的。

考虑 A 的素理想的一条饱和链 \(\mathfrak{p} \subset \mathfrak{p}_1 \subset \ldots \subset \mathfrak{p}_r\)，使得理想 \(\mathfrak{p}_r\) 是极大的。对 \(r\) 作归纳，可设 \(\mathrm{A}_{\mathfrak{p}_1}\)-模 \(\Omega_{\mathfrak{p}_1}\) 是对偶化的。把 \(\Lambda\) 换成 \(\mathrm{A}_{\mathfrak{p}_1}\)，把 \(\mathfrak{p}\) 换成 \(\mathfrak{p}\mathrm{A}_{\mathfrak{p}_1}\)，就化归为环 A 是局部环且链 \(\mathfrak{p} \subset \mathfrak{m}_{\Lambda}\) 为饱和链的情形。

于是设 \(d = \dim(A) = \mathrm{ht}(\mathfrak{m}_{\mathbf{A}})\)。由于 A 是 Macaulay 环（§ 2, 第 2 节，命题 2 的推论），有 \(\dim(A_{\mathfrak{p}}) = \mathrm{ht}(\mathfrak{p}) = d - 1\)。对每个整数 \(i\)，\(A_{\mathfrak{p}}\)-模 \(\operatorname{Ext}_{A_{\mathfrak{p}}}^{i}(\kappa(\mathfrak{p}), \Omega_{\mathfrak{p}})\) 同构于 \(\operatorname{Ext}_{\Lambda}^{i}(A/\mathfrak{p}, \Omega)_{\mathfrak{p}}\)。于是由 § 3, 第 2 节，命题 2，只需证明 \(A/\mathfrak{p}\)-模 \(\operatorname{Ext}_{A}^{i}(A/\mathfrak{p}, \Omega)\) 在 \(i \neq d - 1\) 时为零，而在 \(i = d - 1\) 时秩为 1。

设 \(x\) 为 \(\mathfrak{m}_{\mathrm{A}} - \mathfrak{p}\) 的一个元素，并用 \(\overline{x}\) 表示它在 \(A / \mathfrak{p}\) 中的类。考虑 A-模的正合列

\[
0 \rightarrow \mathrm{A} / \mathfrak {p} \xrightarrow {\overline {{{x}}}} \mathrm{A} / \mathfrak {p} \longrightarrow \mathrm{A} / (\mathfrak {p} + x \mathrm{A}) \rightarrow 0.
\]

A-模 \(\mathrm{A} / (\mathfrak{p} + x\mathrm{A})\) 是有限长度的，因为它的支撑集化为 \(\mathfrak{m}_{\mathrm{A}}\)；由于 A-模 \(\Omega\) 是对偶化的，有 \(\operatorname{Ext}_{\mathrm{A}}^{i}(\mathrm{A} / (\mathfrak{p} + x\mathrm{A}), \Omega) = 0\)（当 \(i \neq d\) 时）（§ 8, 第 5 节，例 3）。于是，由与上述序列以及与 \(\Omega\) 相伴的扩张模正合列推出：以 \(x\) 为比的乘法映射在 A-模 \(\operatorname{Ext}_{\mathrm{A}}^{i}(\mathrm{A} / \mathfrak{p}, \Omega)\) 上是满射（当 \(i \neq d - 1\) 时），这蕴含该模为零（Nakayama 引理）。特别地，\(\operatorname{Ext}_{\mathrm{A}}^{d}(\mathrm{A} / \mathfrak{p}, \Omega)\) 为零，从而得到正合列

\[
0 \rightarrow \operatorname{Ext} _ {\Lambda} ^ {d - 1} (\mathrm{A} / \mathfrak {p}, \Omega) \xrightarrow {x} \operatorname{Ext} _ {\Lambda} ^ {d - 1} (\mathrm{A} / \mathfrak {p}, \Omega) \longrightarrow \operatorname{Ext} _ {\Lambda} ^ {d} (\mathrm{A} / (\mathfrak {p} + x \Lambda), \Omega) \rightarrow 0.
\]

我们有 \(\operatorname{long}_{\mathrm{A}}(\operatorname{Ext}_{\mathrm{A}}^{i}(\mathrm{A}/(\mathfrak{p}+x\mathrm{A}),\Omega))=\operatorname{long}_{\mathrm{A}}(\mathrm{A}/(\mathfrak{p}+x\mathrm{A}))\)（前引）；于是命题由下面的引理得出，把该引理应用于环 B = A/p 及 B-模 \(M=\operatorname{Ext}_{\Lambda}^{d-1}(\mathrm{A}/\mathfrak{p},\Omega)\)：

引理 1. 设 B 为一维的局部 Noether 整环，M 为有限生成的无扭 B-模。假设对 B 的每个非零元素 x，都有 long \(_{B}\) (M/xM) = long \(_{B}\) (B/xB)。那么 B-模 M 的秩为 1。

事实上，设 \(r\) 为 M 的秩；存在 M 的一个秩为 \(r\) 的自由子模 L，使 M/L 为扭模（VII, § 4, 第 1 节，命题 1 的推论），从而是有限长度的（VII, § 2, 第 5 节，引理 1）。M/L 的零化子不化为 0，因而含有 \(x\) 的一个非零元素 \(\mathfrak{m}_{\mathbb{R}}\)。考虑图表

交换图表

\[
\begin{array}{c c c c c c c} 0 \to & \mathrm{L} & \longrightarrow & \mathrm{M} & \longrightarrow & \mathrm{M} / \mathrm{L} & \to 0 \\ & \Biggl \downarrow x _ {\mathrm{L}} & & \Biggl \downarrow x _ {\mathrm{M}} & & \Biggl \downarrow 0 \\ 0 \to & \mathrm{L} & \longrightarrow & \mathrm{M} & \longrightarrow & \mathrm{M} / \mathrm{L} & \to 0. \end{array}
\]

由蛇形引理（A, X, 第 4 页，命题 2），可导出一个正合列

\[
0 \rightarrow \mathrm{M} / \mathrm{L} \longrightarrow \mathrm{L} / x \mathrm{L} \longrightarrow \mathrm{M} / x \mathrm{M} \longrightarrow \mathrm{M} / \mathrm{L} \rightarrow 0,
\]

由此得 long(M/xM) = long(L/xL)。由于按假设 long(M/xM) = long(B/xB)，而 long(L/xL) = r·long(B/xB)，故得 r = 1。

推论 1. 对 A 的每个乘性子集 S，\(\mathrm{S}^{-1}\mathrm{A}\)-模 \(\mathrm{S}^{-1}\Omega\) 都是对偶化的。

推论 2.—— Ω 的支撑集等于 Spec(A)。

事实上，按定义，局部环上的对偶化模是非零的。

推论 3. 设 M 为有限生成的 A-模，i 为整数。A-模 \(\operatorname{Ext}_{\mathrm{A}}^{i}(\mathrm{M},\Omega)\) 是有限生成的，且其支撑集的余维数 \(\geqslant i\)（在 \(\operatorname{Spec}(\mathrm{A})\) 中）。

第一个断言由 A, X, 第 108 页，推论可得。设 \(\mathfrak{p}\) 为 \(\operatorname{Ext}_{\mathrm{A}}^{i}(\mathrm{M},\Omega)\) 的支撑集的一个素理想。有 \(\operatorname{Ext}_{\mathrm{A}}^{i}(\mathrm{M},\Omega)_{\mathfrak{p}} \neq 0\)，从而 \(\operatorname{Ext}_{\mathrm{A}_{\mathfrak{p}}}^{i}(\mathrm{M}_{\mathfrak{p}},\Omega_{\mathfrak{p}}) \neq 0\)（\(\S 3, \mathrm{n}^{\circ}2\)，命题 2），这蕴含 \(\mathrm{di}_{\mathrm{A}_{\mathfrak{p}}}(\Omega_{\mathfrak{p}}) \geqslant i\)。由于 \(\Omega_{\mathfrak{p}}\) 是 \(\mathrm{A}_{\mathfrak{p}}\)-模且是对偶化的（命题 2），有 \(\mathrm{di}_{\mathrm{A}_{\mathfrak{p}}}(\Omega_{\mathfrak{p}}) = \dim(\mathrm{A}_{\mathfrak{p}})\)（命题 1），由此得推论。

命题 3.—— 设 A 为 Noether 局部环，Ω 为对偶化 A-模，而 M 为有限生成的 A-模。

\[
\text { a) } O n a \operatorname{Ext} _ {\mathrm{A}} ^ {i} (\mathrm{M}, \Omega) = 0 \text { for } i <   \dim (\mathrm{A}) \dots \dim_ {\mathrm{A}} (\mathrm{M}).
\]

\begin{enumerate}
\def\labelenumi{\alph{enumi})}
\setcounter{enumi}{1}
\tightlist
\item
  令 \(c = \dim(\mathrm{A}) - \dim_{\mathrm{A}}(\mathrm{M})\)。若 M 非零，则 A-模 \(\operatorname{Ext}_{\mathrm{A}}^{c}(\mathrm{M}, \Omega)\) 不为零。
\end{enumerate}

\[
\mathrm{c)} O n a \operatorname{Ext} _ {\mathrm{A}} ^ {i} (\mathrm{M}, \Omega) = 0 p o u r i > \dim (\mathrm{A}) - \operatorname{prof} _ {\mathrm{A}} (\mathrm{M}).
\]

设 M 非零，并用 F 表示它的支撑集。由 § 1, 第 5 节，命题 9，断言 a) 与 b) 的合取等价于 \(\operatorname{prof}_{\mathbb{F}}(\Omega)=c\)。由于 \(\Omega\) 是 Macaulay 模且其支撑集等于 \(\operatorname{Spec}(\mathbb{A})\)（命题 1 及命题 2 的推论 2），我们有

\[
\operatorname{prof} _ {\mathrm{F}} (\Omega) = \operatorname{codim} (\mathrm{F}, \operatorname{Spec} (\Lambda)) = c
\]

（§ 2, 第 1 节，命题 1 的推论及第 2 节，命题 2 的推论）。

我们对 M 的深度作归纳来证明 c)。若 \(\mathrm{prof}_{\Lambda}(\mathbf{M}) = 0\)，则确实有 \(\mathrm{Ext}_{\Lambda}^{i}(\mathbf{M},\Omega) = 0\)（当 \(i > \dim (\mathbf{A})\) 时），因为 \(\mathrm{di}_{\Lambda}(\Omega) = \dim (\mathbf{A})\)（命题 1）。设 \(\mathrm{prof}_{\Lambda}(\mathbf{M}) > 0\)，则存在 \(x\) 的一个元素 \(\mathfrak{m}_{\mathbf{A}}\)，使以 \(x\) 为比的乘法映射在 M 中是单射的。我们有 \(\mathrm{prof}_{\Lambda}(\mathbf{M} / x\mathbf{M}) = \mathrm{prof}_{\Lambda}(\mathbf{M}) - 1\)（\(\S 1, \mathrm{n}^{\circ}4\)，命题 7）。

考虑扩张模的正合列

\[
\operatorname{Ext} _ {\mathrm{A}} ^ {i} (\mathrm{M}, \Omega) \xrightarrow {x} \operatorname{Ext} _ {\mathrm{A}} ^ {i} (\mathrm{M}, \Omega) \longrightarrow \operatorname{Ext} _ {\Lambda} ^ {i + 1} (\mathrm{M} / x \mathrm{M}, \Omega)
\]

它相伴于正合列

\[
0 \to \mathrm{M} \xrightarrow {x} \mathrm{M} \longrightarrow \mathrm{M} / x \mathrm{M} \to 0.
\]

当 \(i > \dim(\mathrm{A}) - \operatorname{prof}_{\mathrm{A}}(\mathrm{M})\) 时，A-模 \(\operatorname{Ext}_{\mathrm{A}}^{i+1}(\mathrm{M}/x\mathrm{M},\Omega)\) 按归纳假设为零，于是以 \(x\) 为比的乘法映射在 \(\operatorname{Ext}_{\mathrm{A}}^{i}(\mathrm{M},\Omega)\) 上是满射的，这蕴含该 A-模为零（Nakayama 引理）。这就证明了 c)。

推论.—— 若 M 是 Macaulay 模，则当 i ≠ c 时有 Ext \(^{i}\) A(M, Ω) = 0；A-模 Ext \(^{c}\) A(M, Ω) 是 Macaulay 模，且其支撑集等于 M 的支撑集。

第一个断言由命题 3, a) 与 c) 得出。设 \(\mathfrak{p} \in \operatorname{Supp}(M)\)；把 § 2, 第 1 节，命题 1 应用于 M 与 A，我们有

\[
\dim (A _ {p}) - \dim_ {A _ {p}} (M _ {p}) = \dim (A) - \dim_ {A} (M) = c;
\]

由于 \(A_{p}\)-模 \(\Omega_{p}\) 是对偶化的（命题 2），由命题 3, b) 可知 \(A_{p}\)-模 \(\operatorname{Ext}_{A_{p}}^{c}(M_{p},\Omega_{p})\) 不为零。因此，\(\operatorname{Ext}_{A}^{c}(M,\Omega)\) 的支撑集等于 M 的支撑集。

最后，我们对 \(\dim(M)\) 作归纳，证明 A-模 \(\operatorname{Ext}_{A}^{c}(M,\Omega)\) 是 Macaulay 模。当 \(\dim(M)=0\) 时断言成立，因为每个有限长度模都是 Macaulay 模。设 \(\dim(M)>0\)，并选取 \(x\) 的一个元素 \(\mathfrak{m}_{\mathrm{A}}\)，使乘法映射 \(x_{\mathrm{M}}\) 是单射的。A-模 M/\(x\mathrm{M}\) 是 Macaulay 模（\(\S 2\) 第 3 节，命题 4），其维数为 \(\dim(M)-1\)。考虑到上述结果，与正合列 \(0\to\mathrm{M}\xrightarrow{x}\mathrm{M}\longrightarrow\mathrm{M}/x\mathrm{M}\to0\) 相伴的扩张模正合列化为

\[
0 \to \operatorname{Ext} _ {\mathrm{A}} ^ {c} (\mathrm{M}, \Omega) \xrightarrow {x} \operatorname{Ext} _ {\mathrm{A}} ^ {c} (\mathrm{M}, \Omega) \longrightarrow \operatorname{Ext} _ {\mathrm{A}} ^ {c + 1} (\mathrm{M} / x \mathrm{M}, \Omega) \to 0;
\]

于是 § 2, 第 3 节的命题 4 与归纳假设蕴含 Ext \(_{A}^{c}\) (M, Ω) 是 Macaulay 模，由此得推论。

\subsection{对正则序列取商}

命题 4.—— 设 A 为 Noether 环，J 为由 A-正则序列 x 生成的 A 的理想，而 Ω 为有限生成的 A-模。

\begin{enumerate}
\def\labelenumi{\alph{enumi})}
\item
  若 A-模 \(\Omega\) 是对偶化的，则序列 \(\mathbf{x}\) 是 \(\Omega\)-正则的，且 A/J-模 \(\Omega/\mathrm{J}\Omega\) 是对偶化的；
\item
  若 A/J-模 \(\Omega/\mathrm{J}\Omega\) 是对偶化的，J 含于 \(\Lambda\) 的根中，且序列 x 是 \(\Omega\)-正则的，则 A-模 \(\Omega\) 是对偶化的。
\end{enumerate}

对序列 x 的长度作归纳，可化归为它只含单个元素 x 的情形。设 A-模 \(\Omega\) 是对偶化的。对 A 的每个包含 x 的极大理想 m，有 \(\dim(\mathrm{A}_{\mathfrak{m}}/x\mathrm{A}_{\mathfrak{m}})=\dim(\mathrm{A}_{\mathfrak{m}})-1\)（VIII, § 3, 第 1 节，推论 2），从而 \(\mathrm{Hom}_{\mathrm{A}_{\mathfrak{m}}}\left(\mathrm{A}_{\mathfrak{m}} / x\mathrm{A}_{\mathfrak{m}}, \Omega_{\mathfrak{m}}\right) = 0\)（第 1 节，命题 3, a)）。这蕴含 \(\mathrm{Hom}_{\mathrm{A}}(\mathrm{A} / x\mathrm{A}, \Omega) = 0\)，于是乘法映射 \(x_{\Omega}\) 是单射的。因此，为证明命题，可设乘法映射 \(x_{\Omega}\) 是单射的。

令 \(\overline{\mathsf{A}}\) 为环 A/xA；设 \(\mathfrak{m}\) 为 A 的一个包含 \(x\) 的极大理想，并用 \(\overline{\mathfrak{m}}\) 表示它在 \(\overline{\mathsf{A}}\) 中的像。A-模 A/\(\mathfrak{m}\) 被 \(x\) 零化，且等同于 \(\overline{\mathsf{A}}/\overline{\mathfrak{m}}\)；于是对每个整数 \(i \geqslant 1\)，有一个同构 \(\operatorname{Ext}_{\mathsf{A}}^{i}(\mathsf{A}/\mathfrak{m}, \Omega) \longrightarrow \operatorname{Ext}_{\overline{\mathsf{A}}}^{i-1}(\overline{\mathsf{A}}/\overline{\mathfrak{m}}, \Omega/x\Omega)\)（\(\S\) 3, 第 4 节，命题 7）。我们有

\[
\mathrm{ht} (\overline {{\mathfrak {m}}}) = \dim (\overline {{\mathrm{A}}} _ {\overline {{\mathfrak {m}}}}) = \dim (\mathrm{A} _ {\mathfrak {m}} / x \mathrm{A} _ {\mathfrak {m}}) = \dim (\mathrm{A} _ {\mathfrak {m}}) - 1 = \mathrm{ht} (\mathfrak {m}) - 1
\]

（VIII, § 3, 第 1 节，推论 2, a)）。现在，\(\overline{\mathrm{A}}\) 的极大理想正是这样的理想 \(\overline{\mathfrak{m}}\)，其中 \(\mathfrak{m}\) 是 A 的包含 \(x\) 的极大理想；此外，若 \(x\) 属于 A 的根，则 A 的每个极大理想都包含 \(x\)。命题由此得证。

推论 1.—— 设 A 为 Noether 整环。每个对偶化 A-模都是无扭的且秩为 1。

设 \(\Omega\) 为对偶化 A-模；由命题 4，它是无扭的。设 K 为 A 的分式域；K-向量空间 \(K \otimes_{\Lambda} \Omega\) 是对偶化的（\(n^{\circ}\) 1, 命题 2），故维数为 1。

推论 2.—— 设 A 为局部 Macaulay 环，Ω 为有限生成的 A-模，而 x 为 m\(_{A}\) 中元素的一个极大割序列，它生成理想 J。下列条件等价：

\begin{enumerate}
\def\labelenumi{(\roman{enumi})}
\item
  A-模 Ω 是对偶化的；
\item
  A-模 \(\Omega\) 是 Macaulay 模，维数等于 \(\dim(A)\)，且 \(\Omega/J\Omega\) 是局部 Artin 环 \(A/J\) 上的不可分解内射模；
\item
  序列 x 是 Ω-正则的，且 Ω/JΩ 是局部 Artin 环 A/J 上的不可分解内射模；
\item
  序列 \(\mathbf{x}\) 是 \(\Omega\)-正则的，有 \(\mathrm{long}_{\mathrm{A}}(\Omega/\mathrm{J}\Omega) = \mathrm{long}_{\mathrm{A}}(\Lambda/\mathrm{J})\)，且 \(\kappa_{\mathrm{A}}\)-向量空间 \(\mathrm{Hom}_{\Lambda}(\kappa_{\Lambda}, \Omega/\mathrm{J}\Omega)\) 的维数为 1。
\end{enumerate}

(i)⇒(ii)：若 Ω 是对偶化的，则它是 Macaulay 模且维数为 dim(A)（第 1 节，命题 1）。由于 A 是 Macaulay 环，序列 x 是 A-正则的；由命题 4，A/J-模 Ω/JΩ 是对偶化的，从而是 Matlis A/J-模（第 1 节，例 1）。

(ii)⇒(iii)：在假设 (ii) 下，有 \(\dim(\Omega)=\dim(A)\) 与 \(\dim(\Omega/J\Omega)=\dim(A/J)=0\)，故序列 x 对 \(\Omega\) 是割的，从而是 \(\Omega\)-正则的（§ 2, 第 3 节，定理 1）。

(iii)⇒(i)：在 (iii) 的假设下，A/J-模 Ω/JΩ 是 Matlis A/J-模，从而是对偶化的（第 1 节，例 1）；由命题 4，A-模 Ω 是对偶化的。

\begin{enumerate}
\def\labelenumi{(\roman{enumi})}
\setcounter{enumi}{2}
\tightlist
\item
  \(\Leftrightarrow\)（iv）：这由 § 8, 第 3 节的注得出。
\end{enumerate}

\subsection{环的变换}

命题 5.—— 设 \(\rho: A \to B\) 为 Noether 环的同态，它使 B 成为平坦 A-模。假设对 B 的每个极大理想 \(\mathfrak{n}\)，环 \(\kappa(\rho^{-1}(\mathfrak{n})) \otimes_{A} B\) 都是 Gorenstein 环。设 \(\Omega\) 为对偶化 A-模；则 B-模 \(\Omega_{(B)}\) 是对偶化的。

设 \(\mathfrak{n}\) 为 B 的一个极大理想，并用 \(\mathfrak{p}\) 表示它在 A 中的逆像。\(A_{\mathfrak{p}}\)-模 \(B_{\mathfrak{n}}\) 是平坦的，\(A_{\mathfrak{p}}\)-模 \(\Omega_{\mathfrak{p}}\) 是对偶化的，\(\Omega_{(B)} \otimes_B B_{\mathfrak{n}}\) 等同于 \(\Omega_{\mathfrak{p}} \otimes_{A_{\mathfrak{p}}} B_{\mathfrak{n}}\)，而 \(\kappa_{\Lambda_{\mathfrak{p}}} \otimes_{\Lambda_{\mathfrak{p}}} B_{\mathfrak{n}}\)——它等同于 \(\kappa(\mathfrak{p}) \otimes_A B\) 的一个分式环——是 Gorenstein 环。因此，只需在 \(\rho\) 是局部环之间局部同态的情形证明命题，以下即作此假设。

先处理环 A 与 B 都是 Artin 环的情形。令 \(C = B / \mathfrak{m}_{\Lambda}B\)。由于 B 在 A 上平坦，B-模 \(\mathrm{Hom}_{\mathrm{B}}(\mathrm{C},\Omega_{(\mathrm{B})})\) 同构于 \(\mathrm{Hom}_{\Lambda}(\kappa_{\mathrm{A}},\Omega)\otimes_{\mathrm{A}}\mathrm{B}\)（I, § 2, 第 10 节，命题 11），从而同构于 \(\mathrm{Hom}_{\mathrm{A}}(\kappa_{\mathrm{A}},\Omega)\otimes_{\kappa_{\mathrm{A}}}\mathrm{C}\)。由此推出一列同构

\[
\begin{array}{c} \operatorname{Hom} _ {\mathrm{B}} (\kappa_ {\mathrm{B}}, \Omega_ {(\mathrm{B})}) \longrightarrow \operatorname{Hom} _ {\mathrm{C}} (\kappa_ {\mathrm{C}}, \operatorname{Hom} _ {\mathrm{B}} (\mathrm{C}, \Omega_ {(\mathrm{B})})) \longrightarrow \operatorname{Hom} _ {\mathrm{C}} (\kappa_ {\mathrm{C}}, \operatorname{Hom} _ {\mathrm{A}} (\kappa_ {\mathrm{A}}, \Omega) \otimes_ {\kappa_ {\mathrm{A}}} \mathrm{C}) \\ \qquad \qquad \qquad \qquad \qquad \qquad \qquad \qquad \qquad \qquad \longrightarrow \operatorname{Hom} _ {\mathrm{A}} (\kappa_ {\mathrm{A}}, \Omega) \otimes_ {\kappa_ {\mathrm{A}}} \operatorname{Hom} _ {\mathrm{C}} (\kappa_ {\mathrm{C}}, \mathrm{C}). \end{array}
\]

\(\kappa_{\mathrm{A}}\)-向量空间 \(\mathrm{Hom}_{\mathrm{A}}(\kappa_{\mathrm{A}}, \Omega)\) 的维数为 1，因为 \(\Omega\) 是对偶化的；又由于 C 是 Gorenstein 环，\(\kappa_{\mathrm{C}}\)-向量空间 \(\mathrm{Hom}_{\mathrm{C}}(\kappa_{\mathrm{C}}, \mathrm{C})\) 同样如此；于是 \(\kappa_{\mathrm{B}}\)-向量空间 \(\mathrm{Hom}_{\mathrm{B}}(\kappa_{\mathrm{B}}, \Omega_{(\mathrm{B})})\) 的维数为 1。

设 M 为有限长度的 B-模；我们对 \(\operatorname{long}_{\mathrm{B}}(\mathrm{M})\) 作归纳，证明 \(\operatorname{long}_{\mathrm{B}}(\operatorname{Hom}_{\mathrm{B}}(\mathrm{M},\Omega_{(\mathrm{B})})) \leqslant \operatorname{long}_{\mathrm{B}}(\mathrm{M})\)。当 M = 0 时断言显然，而当 \(M = \kappa_{B}\) 时由前所述可得。设 \(\operatorname{long}_{\mathrm{B}}(\mathrm{M}) \geqslant 2\)。存在 B-模的一个正合列

\[
0 \rightarrow \mathrm{M} ^ {\prime} \rightarrow \mathrm{M} \rightarrow \kappa_ {\mathrm{B}} \rightarrow 0
\]

其中 \(\mathrm{long}_{\mathrm{B}}(\mathrm{M}^{\prime}) < \mathrm{long}_{\mathrm{B}}(\mathrm{M})\)。由此导出一个正合列

\[
0 \to \operatorname{Hom} _ {\mathrm{B}} (\kappa_ {\mathrm{B}}, \Omega_ {(\mathrm{B})}) \to \operatorname{Hom} _ {\mathrm{B}} (\mathrm{M}, \Omega_ {(\mathrm{B})}) \to \operatorname{Hom} _ {\mathrm{B}} (\mathrm{M} ^ {\prime}, \Omega_ {(\mathrm{B})}),
\]

然后对 M′ 应用归纳假设即得结论。

设 N 为从 \(\kappa_{A} \otimes_{A} B\) 到 \(\kappa_{B}\) 的典范满射的核。令 \(m = \text{long}_{\text{B}}(\kappa_{\text{A}} \otimes_{\text{A}} \text{B})\)；则有 \(\text{long}_{\text{B}}(\text{N}) = m - 1\)。考虑 B-模的正合列

\[
\begin{array}{c} 0 \to \operatorname{Hom} _ {\mathrm{B}} (\kappa_ {\mathrm{B}}, \Omega_ {(\mathrm{B})}) \longrightarrow \operatorname{Hom} _ {\mathrm{B}} (\kappa_ {\mathrm{A}} \otimes_ {\mathrm{A}} \mathrm{B}, \Omega_ {(\mathrm{B})}) \longrightarrow \operatorname{Hom} _ {\mathrm{B}} (\mathrm{N}, \Omega_ {(\mathrm{B})}) \\ \longrightarrow \operatorname{Ext} _ {\mathrm{B}} ^ {1} (\kappa_ {\mathrm{B}}, \Omega_ {(\mathrm{B})}) \longrightarrow \operatorname{Ext} _ {\mathrm{B}} ^ {1} (\kappa_ {\mathrm{A}} \otimes_ {\mathrm{A}} \mathrm{B}, \Omega_ {(\mathrm{B})}). \end{array}
\]

B-模 \(\mathrm{Hom}_{\mathbb{B}}(\kappa_{\Lambda}\otimes_{\Lambda}\mathrm{B},\Omega_{(\mathbb{B})})\) 与 \(\mathrm{Ext}_{\mathbb{B}}^{1}(\kappa_{\Lambda}\otimes_{\Lambda}\mathrm{B},\Omega_{(\mathbb{B})})\) 分别同构于 \(\mathrm{Hom}_{\Lambda}(\kappa_{\Lambda},\Omega)\otimes_{\Lambda}\mathrm{B}\) 与 \(\mathrm{Ext}_{\mathbb{A}}^{1}(\kappa_{\Lambda},\Omega)\otimes_{\Lambda}\mathrm{B}\)，亦即分别同构于 \(\kappa_{\Lambda}\otimes_{\Lambda}\mathrm{B}\) 与 0。B-模 \(\mathrm{Hom}_{\mathbb{B}}(\kappa_{\mathbb{B}},\Omega_{(\mathbb{B})})\) 与 \(\mathrm{Hom}_{\mathbb{B}}(\kappa_{\mathbb{A}}\otimes_{\mathbb{A}}\mathrm{B},\Omega_{(\mathbb{B})})\) 的长度分别为 1 与 \(m\)，而 \(\mathrm{Hom}_{\mathbb{B}}(\mathrm{N},\Omega_{(\mathbb{B})}))\) 的长度 \(\leqslant m - 1\)；由此推出 B-模 \(\mathrm{Ext}_{\mathbb{B}}^{1}(\kappa_{\mathbb{B}},\Omega_{(\mathbb{B})})\) 为零。由 § 3, 第 3 节，命题 6，B-模 \(\Omega_{(\mathbb{B})}\) 是内射的；从而它是对偶化模（第 1 节，例 1）。

转到一般情形，令 \(C = \kappa_A \otimes_A B\)；按假设它是 Gorenstein 环，因而是 Macaulay 环（\(\S\)。由第 1 节的命题 1，A 是 Macaulay 环，且 A-模 \(\Omega\) 是 Macaulay 模。于是 B 是 Macaulay 环，且 B-模 \(\Omega_{(B)}\) 是 Macaulay 模（\(\S\) 2, 第 7 节，命题 9 的推论 1）。令 \(r = \dim(A)\)，\(s = \dim(C)\)。存在 \(x_1, \ldots, x_r\) 的元素的一个序列（\(\mathfrak{m}_A\)），它对 A-模 A 与 \(\Omega\) 都是正则的；又存在 \(y_1, \ldots, y_s\) 的元素的一个序列（\(\mathfrak{m}_B\)），它对 B-模 C 是正则的；用 \(x\) 表示上述第一个序列在 A 中生成的理想，用 \(y\) 表示上述第二个序列在 B 中生成的理想。序列（\(y_1, \ldots, y_s, \rho(x_1), \ldots, \rho(x_r)\)）对 B-模 B 与 \(\Omega_{(B)}\) 都是正则的（\(\S\) 1, 第 6 节，命题 11），且 A-模 B/\(y\) 是平坦的（前引文，命题 10）。令 \(A' = A/x\)，\(B' = B/(xB + y)\)，并用 \(\rho': A' \to B'\) 表示由 \(\rho\) 通过取商导出的同态。环 \(A'\) 与 \(B'\) 都是 Artin 环，A′-模 \(B'\) 是平坦的，环 \(\kappa_{A'} \otimes_{A'} B'\)——它等同于 C/\(y\)——是 Gorenstein 环（\(\S\) 3, 第 7 节，例 2），且 A′-模 \(\Omega_{(A')}\) 是对偶化的（第 2 节，命题 4）。由证明的第一部分，B′-模 \(\Omega_{(B')}\) 是对偶化的。于是由前引文，B-模 \(\Omega_{(B)}\) 是对偶化的。

推论.—— 设 A 为具有对偶化模 \(\Omega\) 的 Noether 环，B 为 A 上有限多个未定元的多项式代数。则 B-模 \(\Omega_{(B)}\) 是对偶化的。

事实上，对 A 的每个素理想 p，环 \(\kappa(\mathfrak{p})\otimes_{\mathrm{A}}\mathrm{A}[\mathbf{X}]\) 等同于 \(\kappa(\mathfrak{p})[\mathbf{X}]\)，后者是正则的，因而是 Gorenstein 的。

命题 6.—— 设 \(\Lambda\) 为 Noether 局部环，\(\Omega\) 为对偶化 A-模。设 B 为有限 A-代数；假设 A-模 B 是 Macaulay 模。B-模 \(\mathrm{Ext}_{\mathrm{A}}^{i}(\mathrm{B},\Omega)\) 当 \(i\neq \dim (\Lambda) - \dim (\mathrm{B})\) 时为零，而当 \(i = \dim (\Lambda) - \dim (\mathrm{B})\) 时是对偶化的。

我们有 \(\dim(B) = \dim_{A}(B) \leqslant \dim(\Lambda)\)（VIII, § 2, 第 3 节，定理 1 c)）；令 \(c = \dim(A) - \dim(B)\)。有 \(\operatorname{Ext}_{A}^{i}(B, \Omega) = 0\)（其中 \(i \neq c\)），因为 A-模 B 是 Macaulay 模（第 1 节，命题 3 的推论）。我们来证明 B-模 \(\operatorname{Ext}_{\Lambda}^{c}(B, \Omega)\) 是对偶化的。

先设 \(\dim(B) = 0\)。B 的谱 X 是有限的，且由极大理想组成（IV, § 2, 第 5 节，命题 9）；典范映射 \(B \to \prod_{n \in X} B_n\) 是同构（前引文，推论 1）。于是 B-模 \(\Omega' = \text{Ext}_A^c(B, \Omega)\) 是诸模 \(\text{Ext}_A^c(B_n, \Omega)\) 的直和；由于 \(\text{Ext}_A^c(B_n, \Omega)\) 的支撑集含于 \(\{\mathfrak{n}\}\)，它等同于 \(\Omega'_n\)。有 \(\dim(B_n) = 0\)（对每个 \(\mathfrak{n}\)）；因此，为证明 B-模 \(\Omega'\) 是对偶化的，只需证明 \(B_n\)-模 \(\text{Ext}_A^c(B_n, \Omega)\)（对每个 \(\mathfrak{n} \in X\)）都如此，这就化归为环 B 是局部环的情形。在此情形，由 § 8, 第 5 节，例 6，B-模 \(\text{Ext}_A^c(B, \Omega)\) 同构于 \(\text{Hom}_A(B, I)\)，其中 I 是 Matlis A-模；从而它是 Matlis B-模（§ 8, 第 6 节，命题 5 的推论），因此是对偶化 B-模（第 1 节，例 1）。

现设 \(\dim(B) > 0\)，并对 \(\dim(B)\) 作归纳推理。有 \(\operatorname{prof}_{A}(B) = \dim_{A}(B) = \dim(B)\)，从而 \(\operatorname{prof}_{A}(B) > 0\)；另一方面，有 \(\operatorname{prof}(A) = \dim(A) > 0\)（第 1 节，命题 1），于是 \(\operatorname{prof}_{A}(A \oplus B) > 0\)。因此存在 \(x\) 的一个元素 \(\mathfrak{m}_{\Lambda}\)，使同态 \(x_{\Lambda}\) 与 \(x_{B}\) 都是单射的。

考虑与正合列 \(0 \to \mathrm{B} \xrightarrow{x_{\mathrm{B}}} \mathrm{B} \longrightarrow \mathrm{B}/x\mathrm{B} \to 0\) 以及与 A-模 \(\Omega\) 相伴的扩张模正合列。A-模 B/xB 是 Macaulay 模（\(\S 2, n^{\circ} 1\)，例 3），维数为 \(\dim(\mathrm{B}) - 1\)（VIII, \(\S 3, n^{\circ} 2\)）；因此有 \(\operatorname{Ext}_{\mathrm{A}}^{i}(\mathrm{B}/x\mathrm{B}, \Omega) = 0\)（对 \(i \neq c + 1\)）（\(n^{\circ} 1\)，命题 3 的推论）。由于有 \(\operatorname{Ext}_{\mathrm{A}}^{i}(\mathrm{B}, \Omega) = 0\)（对 \(i \neq c\)），故得到 B-模的一个正合列

\[
0 \rightarrow \operatorname{Ext} _ {\mathrm{A}} ^ {c} (\mathrm{B}, \Omega) \xrightarrow {x} \operatorname{Ext} _ {\mathrm{A}} ^ {c} (\mathrm{B}, \Omega) \longrightarrow \operatorname{Ext} _ {\mathrm{A}} ^ {c + 1} (\mathrm{B} / x \mathrm{B}, \Omega) \rightarrow 0.
\]

由归纳假设，B/xB-模 \(\operatorname{Ext}_{\Lambda}^{c+1}(\mathrm{B}/x\mathrm{B},\Omega)\) 是对偶化的。由于 A-代数 B 是有限的，\(m_{A}\) 在 B 中的像含于 B 的根（V, § 2, 第 1 节，命题 1）；由第 2 节的命题 4，B-模 \(\operatorname{Ext}_{\mathrm{A}}^{c}(\mathrm{B},\Omega)\) 是对偶化的。

推论 1.—— 设 A 为 Noether 环，Ω 为对偶化 A-模，B 为有限 A-代数；假设 A-模 B 是 Macaulay 模。则 B-模 Ext \(_{A}\) (B, Ω) 是对偶化的。

设 \(\Omega'\) 为 B-模 \(\mathrm{Ext}_{\mathsf{A}}(\mathsf{B},\Omega)\)。设 \(\mathfrak{n}\) 为 B 的一个极大理想；它在 A 中的逆像是极大理想 \(\mathfrak{m}\)（V, § 2, 第 1 节，命题 1）。\(\Lambda_{\mathfrak{m}}\)-代数 \(\mathrm{B}_{\mathfrak{m}} = \mathrm{A}_{\mathfrak{m}} \otimes_{\mathsf{A}} \mathrm{B}\) 是有限的，且是 Macaulay \(\mathrm{A}_{\mathfrak{m}}\)-模；由命题，\(\mathrm{B}_{\mathfrak{m}}\)-模 \(\Omega_{\mathfrak{m}}'\)——它等同于 \(\mathrm{Ext}_{\mathrm{A}_{\mathfrak{m}}}(B_{\mathfrak{m}},\Omega_{\mathfrak{m}})\)（§ 3, 第 2 节，命题 2）——是对偶化的。由于 \(\mathrm{B}_{\mathfrak{n}}\) 是 \(\mathrm{B}_{\mathfrak{m}}\) 的一个分式环，\(\mathrm{B}_{\mathfrak{n}}\)-模 \(\Omega_{\mathfrak{n}}'\) 是对偶化的，推论得证。

注.—— 保留推论 1 的假设，并进一步设典范同态 \(\rho : A \to B\) 是单射的。于是有 \(\dim(\Lambda_{\mathfrak{m}}) = \dim(B_{\mathfrak{m}})\)，其中 \(\mathfrak{m}\) 取遍 \(A\) 的极大理想。由命题 6 与推论 1，\(\operatorname{Ext}_{A}^{i}(B, \Omega)\) 当 \(i \neq 0\) 时为零，且 B-模 \(\operatorname{Hom}_{A}(B, \Omega)\) 是对偶化的。

推论 2.—— 若 Noether 环 A 具有对偶化模，则每个是 Macaulay 环的有限生成 A-代数都具有对偶化模。

这由推论 1 及命题 5 的推论得出。

推论 3.—— 每个可表现的 Macaulay 环（特别地，每个完备局部 Macaulay 环）都具有对偶化模。

事实上，设 R 为正则环，A 为 R 的一个 Macaulay 商环。R-模 A 是 Macaulay 模（§ 2, 第 5 节，例 5），且 R 具有对偶化模（第 1 节，例 2）；于是由推论 1，A 亦然。此外，已经指出完备局部 Noether 环是可表现的（§ 4, 第 4 节，命题 6, c)）。

更一般地，每个是 Gorenstein 环之商的 Macaulay 环都具有对偶化模。反之，可以证明具有对偶化模的局部 Macaulay 环是局部 Gorenstein 环的商（习题 1）。

\subsection{对偶化模的结构}

引理 2.—— 设 A 为 Noether 环，M 与 N 为有限生成的 A-模，\(u: \mathrm{M} \to \mathrm{N}\) 为一个同态。设 \(x\) 为 A 的根的一个元素，使同态 \(x_{\mathrm{N}}\) 是单射的。若同态 \(\overline{u}: \mathrm{M}/x\mathrm{M} \to \mathrm{N}/x\mathrm{N}\)（它由 \(u\) 诱导）是单射的（相应地，满射的、双射的），则 \(u\) 亦然。

关于 \(u\) 的满射性的断言由 Nakayama 引理（II, § 3, 第 2 节，命题 4 的推论 1）得出，无须对 \(x_{\mathrm{N}}\) 作任何假设。考虑行正合的交换图表

\[
\begin{array}{c c c c c c c} & \mathrm{M} & \xrightarrow {x _ {\mathrm{M}}} & \mathrm{M} & \longrightarrow & \mathrm{M} / x \mathrm{M} \\ & u \Bigg \downarrow & & u \Bigg \downarrow & & \overline {{u}} \Bigg \downarrow \\ 0 & \longrightarrow & \mathrm{N} & \xrightarrow {x _ {\mathrm{N}}} & \mathrm{N} & \longrightarrow & \mathrm{N} / x \mathrm{N} \end{array} ;
\]

利用蛇形引理（I, § 1, 第 4 节，命题 2），可导出一个正合列 \(\operatorname{Ker} u \xrightarrow{x} \operatorname{Ker} u \longrightarrow \operatorname{Ker} \overline{u}\)。若 \(\overline{u}\) 是单射的，则以 \(x\) 为比的乘法映射在 \(\operatorname{Ker} u\) 上是满射的，由 Nakayama 引理这蕴含 \(\operatorname{Ker} u = 0\)。

命题 7.—— 设 A 为 Noether 环，Ω 为对偶化 A-模。

\begin{enumerate}
\def\labelenumi{\alph{enumi})}
\item
  有 \(\operatorname{Ext}_{\mathrm{A}}^{i}(\Omega, \Omega) = 0\)（对 \(i > 0\)）。
\item
  \(L'\) 典范同态 \(\gamma : \mathrm{A} \to \operatorname{End}_{\mathrm{A}}(\Omega)\) 是双射的。
\item
  每个对偶化 A-模都具有形状 \(\Omega\otimes_{\mathrm{A}}\mathrm{L}\)，其中 \(\mathrm{L}\) 是秩为 1 的投射 A-模。
\end{enumerate}

\begin{enumerate}
\def\labelenumi{\Alph{enumi})}
\tightlist
\item
  先处理环 A 是局部环的情形。在此情形，条件 c) 只是说两个对偶化模是同构的。
\end{enumerate}

设 \(\Omega'\) 为对偶化 A-模。有 \(\operatorname{prof}_{\mathrm{A}}(\Omega') = \dim_{\mathrm{A}}(\Omega') = \dim(\mathrm{A})\)（第 1 节，命题 1），故有 \(\operatorname{Ext}_{\Lambda}^{i}(\Omega', \Omega) = 0\)（对 \(i \neq 0\)）（第 1 节，命题 3, c)），由此得 a)。

我们对整数 dim(A)（它等于 prof(A)）作归纳，来证明 b) 与 c)。若它为零，则环 A 是 Artin 环，\(\Omega'\) 与 \(\Omega\) 都是 Matlis A-模（\(n^{\circ}\) 1, 例 1）；因而它们是同构的（\(\S\) 8, \(n^{\circ}\)），且典范映射 A → \(\mathrm{End}_{\mathrm{A}}(\Omega)\) 是双射的（\(\S\) 8, \(n^{\circ}\) 2, 命题 3, c)）。设 dim(A) \textgreater{} 0，并设 \(x\) 为 \(\mathfrak{m}_{\mathrm{A}}\) 的一个可消去元素。同态 \(x_{\Omega}\) 是单射的（\(n^{\circ}\) 2, 命题 4），并且有一个正合列

\[
0 \rightarrow \Omega \xrightarrow {x _ {\Omega}} \Omega \longrightarrow \Omega / x \Omega \rightarrow 0.
\]

由于 \(\mathrm{Ext}_{\mathrm{A}}^{1}(\Omega',\Omega)\) 为零，且 \(\mathrm{Hom}_{\mathrm{A}}(\Omega',\Omega /x\Omega)\) 等同于 \(\mathrm{Hom}_{\mathrm{A} / x\mathrm{A}}(\Omega'/x\Omega',\Omega/x\Omega)\)，可导出一个正合列

\[
0 \rightarrow \operatorname{Hom} _ {\mathrm{A}} \left(\Omega^ {\prime}, \Omega\right) \xrightarrow {x} \operatorname{Hom} _ {\mathrm{A}} \left(\Omega^ {\prime}, \Omega\right) \xrightarrow {p} \operatorname{Hom} _ {\mathrm{A} / x \mathrm{A}} \left(\Omega^ {\prime} / x \Omega^ {\prime}, \Omega / x \Omega\right)\rightarrow 0,\tag{1}
\]

其中 \(p\) 是典范映射。由命题 4，A/xA-模 \(\Omega/x\Omega\) 与 \(\Omega'/x\Omega'\) 都是对偶化的，因而由归纳假设它们同构。设 \(\overline{u}\) 是从 \(\Omega'/x\Omega'\) 到 \(\Omega/x\Omega\) 的一个同构。考虑到正合列 (1)，存在一个 A-同态 \(u: \Omega' \to \Omega\)，使 \(p(u) = \overline{u}\)；由引理 2，\(u\) 是双射的，这就证明了 c)。由归纳假设，典范同态 A/xA \(\longrightarrow\) End \(_{A/xA}(\Omega/x\Omega)\) 是双射的。鉴于正合列 (1)，此同态等同于同态 \(\overline{\gamma}: A/xA \longrightarrow \text{End}_A(\Omega)/x \text{End}_A(\Omega)\)（它由 \(\gamma\) 诱导）；于是由引理 2 推出 \(\gamma\) 是双射的，由此得 b)。

\begin{enumerate}
\def\labelenumi{\Alph{enumi})}
\setcounter{enumi}{1}
\tightlist
\item
  转到一般情形。对 A 的每个极大理想 m 及每个整数 i \textgreater{} 0，由前所述有 \(\operatorname{Ext}_{A_{m}}^{i}(\Omega_{\mathfrak{m}}, \Omega_{\mathfrak{m}}) = 0\)，因此 \(\operatorname{Ext}_{A}^{i}(\Omega, \Omega)_{\mathfrak{m}} = 0\)（§ 3, 第 2 节，命题 2），这蕴含 \(\operatorname{Ext}_{A}^{i}(\Omega, \Omega) = 0\)（II, § 3, 第 3 节，定理 1 的推论 2）。同样，同态 \(\gamma_{m}: A_{m} \to \operatorname{End}_{A}(\Omega)_{m}\) 对 A 的每个极大理想 m 都是双射的，故 \(\gamma\) 是双射的（前引文，定理 1）。
\end{enumerate}

最后证明 c)。设 \(\Omega'\) 为对偶化 A-模。用 L 表示 A-模 \(\mathrm{Hom}_{\mathrm{A}}(\Omega',\Omega)\)，并用 \(v:\Omega' \otimes_{\mathrm{A}} \mathrm{L} \longrightarrow \Omega\) 表示满足 \(v(x \otimes f) = f(x)\)（对 \(x \in \Omega'\)，\(f \in \mathrm{L}\)）的同态。设 \(\mathfrak{m}\) 为 A 的一个极大理想。\(\mathrm{A}_{\mathfrak{m}}\)-模 \(\mathrm{L}_{\mathfrak{m}}\) 等同于 \(\mathrm{Hom}_{\mathrm{A}_{\mathfrak{m}}}(\Omega_{\mathfrak{m}}',\Omega_{\mathfrak{m}})\)；由已处理的情形，它是秩为 1 的自由模，且每个同构 \(h:\Omega_{\mathfrak{m}}' \to \Omega_{\mathfrak{m}}\) 都是它的一个生成元。当利用生成元 \(\mathrm{L}_{\mathfrak{m}}\) 把 \(\mathrm{A}_{\mathfrak{m}}\) 与 \(h\) 等同时，同态 \(v_{\mathfrak{m}}:\Omega_{\mathfrak{m}}' \otimes_{\mathrm{A}_{\mathfrak{m}}}\mathrm{L}_{\mathfrak{m}} \longrightarrow \Omega_{\mathfrak{m}}\) 等同于 \(h\)，因此它是双射的。由于这对 A 的每个极大理想 \(\mathfrak{m}\) 都成立，A-模 L 是秩为 1 的投射模（II, § 5, 第 3 节，定理 2），且同态 \(v\) 是双射的（II, § 3, 第 3 节，定理 1）。

推论 1.—— 为使 A 是 Gorenstein 环，必须且只须 A-模 Ω 是秩为 1 的投射模。

这由第 1 节的例 2 及命题 7 c) 得出。

推论 2.—— 设环 A 是可表现的。A 的使 A\(_{p}\) 为 Gorenstein 环的素理想 p 的集合在 Spec(A) 中是开的。

设 p 为 A 的一个使 \(A_{p}\) 是 Gorenstein 环的素理想。于是它是 Macaulay 环；必要时把 A 换成 \(A_{f}\)，其中 f 是 A − p 的一个适当元素，就可化归为 A 是 Macaulay 环的情形（§ 4, 第 4 节，命题 7, c)）。设 \(\Omega\) 为对偶化 A-模（第 3 节，命题 6 的推论 3）。于是 \(\Omega_{p}\) 是 Gorenstein 环 \(A_{p}\) 上的对偶化模（第 1 节，命题 2），从而是秩为 1 的自由模（推论 1）。因此存在 A − p 的一个元素 g，使 \(A_{g}\)-模 \(\Omega_{g}\) 是秩为 1 的自由模（II, § 5, 第 1 节，命题 2 的推论）。于是 \(A_{g}\) 是 Gorenstein 环（推论 1），且对 A 的每个不含 g 的素理想 q，\(A_{q}\) 亦然（§ 3, 第 7 节，例 1），这就证明了推论。

\subsection{有限生成模的对偶性}

在本节中，我们考虑一个具有对偶化模 \(\Omega\) 的有限维 Noether 环 A。于是有 \(\mathrm{di}_{\Lambda}(\Omega) = \dim(\mathrm{A}) < +\infty\)（\(n^{\circ}\) 1, 命题 1）。选取一个有限长度的内射分解 \(e: \Omega \to (\mathbf{I}, \delta)\)。对 A-模的每个复形 C，用 \(\mathbf{D}(\mathbf{C})\) 表示复形 \(\operatorname{Homgr}_{\Lambda}(\mathbf{C}, \mathbf{I})\)。这特别适用于每个 A-模 M，把它视为集中在 0 次的复形；于是有 \(\mathbf{D}(\mathrm{M})^{i} = \mathrm{Hom}_{\mathrm{A}}(\mathrm{M},\mathrm{I}^{i})\)（对每个整数 \(i\)）。回忆在 A, X, 第 100 页，定理 1 中，我们构造了一个典范同构

\[
\varphi (\mathrm{M}, \mathrm{I}): \mathrm{H} (\mathbf {D} (\mathrm{M})) \longrightarrow \operatorname{Ext} _ {\mathrm{A}} (\mathrm{M}, \Omega).
\]

例.—— 1) 复形 \(\mathbf{D}(\mathbf{A}) = \operatorname{Homgr}_{\mathbf{A}}(\mathbf{A},\mathbf{l})\) 等同于 I。映射 \(c:\Omega \to \mathbf{D}(\Lambda)\) 按定义是一个拟同构。

\begin{enumerate}
\def\labelenumi{\arabic{enumi})}
\setcounter{enumi}{1}
\item
  同态 \(e \in \text{Homgr}_A(\Omega, I)^0\) 是 \(\mathbf{D}(\Omega)^0\) 的一个元素；A-线性映射 \(\tilde{e}: A \to \mathbf{D}(\Omega)\)（满足 \(\tilde{e}(1) = e\)）是一个拟同构（第 4 节，命题 7, a) 与 b)）。
\item
  设 S 为 A 的一个乘性子集。\(S^{-1}A\)-模 \(S^{-1}\Omega\) 是对偶化的（第 1 节，命题 2 的推论 1）；\(S^{-1}A\)-模 \(S^{-1}I^{i}\) 都是内射的（§ 3, 第 2 节，命题 3 的推论 1），而由 e 导出的态射 \(e^{\prime}:S^{-1}\Omega\to S^{-1}I\) 是 \(S^{-1}\Omega\) 的一个内射分解，因此可以对它应用前面的结果。对每个有限型的复形 C（特别地，对每个有限生成的 A-模 M），从 \(S^{-1}\mathbf{D}(\mathbf{C})\) 到 \(\operatorname{Homgr}_{\mathrm{S}^{-1}\mathrm{A}}(\mathrm{S}^{-1}\mathrm{C},\mathrm{S}^{-1}\mathrm{I})=\mathbf{D}(\mathrm{S}^{-1}\mathrm{C})\) 的典范同态是双射的。
\item
  设 A 为 Dedekind 环，K 为它的分式域。A-模 A 是对偶化的，且具有由正合列定义的长度为 1 的内射分解 I
\end{enumerate}

\[
0 \to \mathrm{A} \stackrel {e} {\longrightarrow} \mathrm{K} \stackrel {\delta} {\longrightarrow} \mathrm{K/A} \to 0
\]

其中 δ 是典范满射。对每个 A-模 M，复形 D(M) 是集中在次数 0 与 1 的复形

\[
\dots \longrightarrow 0 \longrightarrow \operatorname{Hom} _ {\mathrm{A}} (\mathrm{M}, \mathrm{K}) \stackrel {d} {\longrightarrow} \operatorname{Hom} _ {\mathrm{A}} (\mathrm{M}, \mathrm{K} / \mathrm{A}) \longrightarrow 0 \longrightarrow \dots
\]

其中 \(d = \mathrm{Hom}_{\mathrm{A}}(1_{\mathrm{M}},\delta)\)。我们有一个正合列

\[
0 \to \operatorname{Hom} _ {\mathrm{A}} (\mathrm{M}, \mathrm{A}) \longrightarrow \mathbf {D} (\mathrm{M}) ^ {0} \stackrel {d} {\longrightarrow} \mathbf {D} (\mathrm{M}) ^ {1} \longrightarrow \operatorname{Ext} _ {\mathrm{A}} ^ {1} (\mathrm{M}, \mathrm{A}) \to 0.
\]

对复形的每个态射 \(f: \mathrm{C} \to \mathrm{C}'\)，用 \(\mathbf{D}(f): \mathbf{D}(\mathrm{C}') \to \mathbf{D}(\mathrm{C})\) 表示复形的态射 \(\operatorname{Homgr}_{\mathrm{A}}(f, 1_{\mathrm{I}})\)。若 \(f\) 是拟同构，则 \(\mathbf{D}(f)\) 是拟同构（A, X, 第 86 页，命题 4, b)）。若 \(\mathrm{C}' \xrightarrow{f} \mathrm{C} \xrightarrow{g} \mathrm{C}''\) 是复形的一个正合列，则复形列 \(\mathbf{D}(\mathrm{C}''') \xrightarrow{\mathbf{D}(g)} \mathbf{D}(\mathrm{C}) \xrightarrow{\mathbf{D}(f)} \mathbf{D}(\mathrm{C}')\) 是正合的（A, X, 第 83 页，命题 2, a)）。

设 M 为 A-模。对 M 的每个元素 m，结合一个从 D(M) 到 I 的映射 \(\alpha_{\mathrm{M}}(m): f \mapsto f(m)\)；它是 \(\mathbf{D}(\mathbf{D}(\mathbf{M}))_{0} = \operatorname{Homgr}_{\mathrm{A}}(\mathbf{D}(\mathbf{M}), \mathrm{I})_{0}\) 的一个元素。由定义可知，\(\alpha_{\mathrm{M}}(m)\) 是复形的一个态射，从而是 \(\mathrm{Z}_{0}(\mathbf{D}(\mathbf{D}(\mathbf{M})))\) 的一个元素。

这样就定义了复形的一个态射：

\[
\alpha_ {\mathrm{M}}: \mathrm{M} \rightarrow \mathbf {D} (\mathbf {D} (\mathrm{M})) ,
\]

由此，通过取同调，得到 A-模的一个同态

\[
\alpha_ {\mathrm{M}}: \mathrm{M} \to \mathrm{H} _ {0} (\mathbf {D} (\mathbf {D} (\mathrm{M}))).
\]

定理 1. 设 M 为有限生成的 A-模。则 \(\alpha_{\mathrm{M}}\) 是一个拟同构：有 \(\mathrm{H}_{i}(\mathbf{D}(\mathbf{D}(\mathbf{M}))) = 0\)（对 \(i \neq 0\)），且同态 \(\alpha_{\mathrm{M}}\) 是双射的。

先取 M = A。映射 \(e : \Omega \to \mathbf{D}(A)\) 是一个拟同构（例 1），故映射 \(\mathbf{D}(e) : \mathbf{D}(\mathbf{D}(A)) \to \mathbf{D}(\Omega)\) 亦然。映射 \(\tilde{e} : A \to \mathbf{D}(\Omega)\) 是一个拟同构（例 2），且有 \(\mathbf{D}(e) \circ \alpha_{\mathbf{A}} = \tilde{c}\)；于是 \(\alpha_{A}\) 是拟同构，这就证明了此情形下的定理。由此推出，当 A-模 M 是有限生成的自由模时，\(\alpha_{M}\) 是拟同构。

转到一般情形；我们将对整数 n 作归纳，证明下列断言：

\((A_{n})\) 对每个有限生成的 A-模 M，同态 \(H_{i}(\alpha_{M})\) 当 \(i \leqslant n\) 时是双射的。

这也意味着 \(\mathrm{H}_{i}(\mathbf{D}(\mathbf{D}(\mathbf{M})))\) 在 \(i \neq 0\) 且 \(i \leqslant n\) 时为零，且 \(\alpha_{M}\) 当 \(n \geqslant 0\) 时是双射的。注意，\((\mathrm{A}_{n})\) 当 n \textless{} -d 时成立，其中 d 是复形 I 的长度：事实上，A-模 \(\mathbf{D}(\mathbf{D}(\mathbf{M}))_{i}\) 等于 \(\bigoplus_{p}\mathrm{Hom}_{\mathrm{A}}(\mathrm{Hom}_{\mathrm{A}}(\mathrm{M},\mathrm{I}^{p}),\mathrm{I}^{p-i})\)，故当 i \textless{} -d 及 i \textgreater{} d 时它为零。

我们来证明蕴含关系 \((\mathrm{A}_{n}) \Rightarrow (\mathrm{A}_{n+1})\)。设 M 为有限生成的 A-模。存在一个有限生成的自由 A-模 L 及一个正合列 \(0 \to N \xrightarrow{u} L \xrightarrow{v} M \to 0\)。序列 \(0 \to \mathbf{D}(M) \xrightarrow{\mathbf{D}(v)} \mathbf{D}(L) \xrightarrow{\mathbf{D}(u)} \mathbf{D}(N) \to 0\) 是正合的；同样，若令 \(u' = \mathbf{D}(\mathbf{D}(u))\) 与 \(v' = \mathbf{D}(\mathbf{D}(v))\)，则序列 \(0 \to \mathbf{D}(\mathbf{D}(N)) \xrightarrow{u'} \mathbf{D}(\mathbf{D}(L)) \xrightarrow{v'} \mathbf{D}(\mathbf{D}(M)) \to 0\) 是正合的。

由于 \(H_{i}(\mathbf{D}(\mathbf{D}(\mathrm{L})))\)（对 \(i \neq 0\)）为零，我们有同构

\[
\mathrm{H} _ {i} (\mathbf {D} (\mathbf {D} (\mathrm{M}))) \longrightarrow \mathrm{H} _ {i - 1} (\mathbf {D} (\mathbf {D} (\mathrm{N}))) \quad \text { pour } i \neq 0, 1;
\]

这就蕴含了蕴含关系 \((\mathrm{A}_{n})\Rightarrow(\mathrm{A}_{n+1})\)（对 \(n\neq-1\) 与 \(n\neq0\)）。考虑行正合的交换图表

\[
\begin{array}{c c c c c c c} 0 \to & N & \xrightarrow {u} & L & \xrightarrow {v} & M & \to 0 \\ & \alpha_ {N} \Bigg \downarrow & & \alpha_ {L} \Bigg \downarrow & & \alpha_ {M} \Bigg \downarrow \\ 0 \to H _ {1} (\mathbf {D} (\mathbf {D} (M))) & \xrightarrow {} H _ {0} (\mathbf {D} (\mathbf {D} (N))) & \xrightarrow {H _ {0} (u ^ {\prime})} & H _ {0} (\mathbf {D} (\mathbf {D} (L))) & \xrightarrow {H _ {0} (v ^ {\prime})} & H _ {0} (\mathbf {D} (\mathbf {D} (M))) \longrightarrow H _ {- 1} (\mathbf {D} (\mathbf {D} (N)) \end{array}
\]

其中 \(\alpha_{\mathrm{L}}\) 是双射的。若 \((\mathrm{A}_0)\) 成立，则同态 \(\alpha_{\mathrm{N}}\) 也是双射的，故 \(\mathrm{H}_0(u')\) 是单射的，且得到 \(\mathrm{H}_1(\mathbf{D}(\mathbf{D}(\mathbf{M}))) = 0\)，由此得 \((\mathrm{A}_1)\)。若 \((\mathrm{A}_{-1})\)

成立，则 \(\mathrm{H}_{-1}(\mathbf{D}(\mathbf{D}(\mathrm{N})))\) 为零，故 \(\mathrm{H}_{0}(v^{\prime})\) 是满射的，这蕴含 \(\alpha_{M}\) 是满射的。由于这对每个有限生成的 A-模 M 都成立，\(\alpha_{N}\) 也是满射的；由 I, § 1, 第 4 节，命题 2 的推论 2，\(\alpha_{M}\) 是双射的，于是 \((\mathrm{A}_{0})\) 成立。

于是 \((\mathrm{A}_{n})\) 对所有 n 都成立，这就证明了定理。

设 M 为有限生成的 A-模；令 \(c = \dim(A) - \dim_{A}(M)\)。令 \(\mathbf{D}'(M)\) 为 \(\mathbf{D}(M)\) 的下述子复形：它等于 \(\bigoplus_{i\in\mathcal{O}}\mathbf{D}(M)^{i}\bigoplus Z^{c}(\mathbf{D}(M))\)；并用

\[
j: \mathbf {D} ^ {\prime} (\mathrm{M}) \rightarrow \mathbf {D} (\mathrm{M})
\]

表示典范单射。由典范满射 \(\mathrm{Z}^{c}(\mathbf{D}(\mathrm{M}))\to\mathrm{H}^{c}(\mathbf{D}(\mathrm{M}))\) 与同构 \(\varphi(\mathrm{M},\mathrm{I})\) 导出复形的一个态射

\[
p: \mathbf {D} ^ {\prime} (\mathrm{M}) (- c) \rightarrow \operatorname{Ext} _ {\Lambda} ^ {c} (\mathrm{M}, \Omega);
\]

由于 \(\mathrm{H}^{i}(\mathbf{D}(\mathrm{M}))\) 当 i \textless{} c 时为零（第 1 节，命题 3 a)），\((\mathbf{D}^{\prime}(\mathrm{M})(-c), p)\) 是 \(\operatorname{Ext}_{\mathrm{A}}^{c}(\mathrm{M}, \Omega)\) 的一个左分解。由 A, X, 第 100 页，定理 1，我们有一个典范同构

\[
\varphi^ {0} \left(\mathbf {D} ^ {\prime} (\mathrm{M}) (- c), \mathrm{I}\right): \mathrm{H} _ {0} \left(\mathbf {D} \left(\mathbf {D} ^ {\prime} (\mathrm{M})\right)\right) \longrightarrow \operatorname{Ext} _ {\mathrm{A}} ^ {c} \left(\operatorname{Ext} _ {\mathrm{A}} ^ {c} (\mathrm{M}, \Omega), \Omega\right);
\]

把它与同态 \(\mathrm{H}_{0}(\mathbf{D}(j)) : \mathrm{H}_{0}(\mathbf{D}(\mathbf{D}(\mathrm{M}))) \to \mathrm{H}_{0}(\mathbf{D}(\mathbf{D}'(\mathrm{M})))\) 复合，便得到一个同态

\[
\mathrm{H} _ {0} (\mathbf {D} (\mathbf {D} (\mathrm{M}))) \longrightarrow \operatorname{Ext} _ {\mathrm{A}} ^ {c} \left(\operatorname{Ext} _ {\mathrm{A}} ^ {c} (\mathrm{M}, \Omega), \Omega\right),
\]

最后，通过与 \(\alpha_{\mathrm{M}}\) 复合，便得到一个典范同态

\[
\beta_ {\mathrm{M}}: \mathrm{M} \longrightarrow \operatorname{Ext} _ {\mathrm{A}} ^ {c} (\operatorname{Ext} _ {\mathrm{A}} ^ {c} (\mathrm{M}, \Omega), \Omega).
\]

推论.—— 若 A-模 M 是 Macaulay 模，则同态 \(\beta_{\mathrm{M}}\) 是双射的。

若 M 是 Macaulay 模，则 \(\Lambda\)-模 \(\mathrm{H}^i(\mathbf{D}(\mathrm{M}))\) 当 \(i \neq c\) 时为零（\(n^\circ 1\)，命题 3 的推论），故典范单射 \(j: \mathbf{D}'(\mathrm{M}) \to \mathbf{D}(\mathrm{M})\) 是拟同构；于是复形的态射 \(\mathbf{D}(j): \mathbf{D}(\mathbf{D}(\mathrm{M})) \to \mathbf{D}(\mathbf{D}'(\mathrm{M}))\) 是拟同构（\(\Lambda\), X, 第 86 页，命题 4）。从而 \(\mathrm{H}_0(\mathbf{D}(j))\) 是双射的；另一方面，由定理 1，\(\alpha_{\mathrm{M}}\) 是双射的，推论得证。

当 A-模 M 是有限长度时，A-模 \(\operatorname{Ext}_{\Lambda}^{c}(M,\Omega)\) 等同于 M 的 Matlis 对偶（参看 § 8, 第 5 节，例 3 及定理 3），由此重新得到 § 8, 第 4 节的命题 4。

\subsection{例子：一维的情形}

在本节中，考虑一个一维的 Noether 整环 A，它具有对偶化模 \(\Omega\)。用 K 表示 A 的分式域，用 V 表示 K-向量空间 \(K \otimes_{A} \Omega\)。

典范同态 \(\Omega\to V\) 是单射的，且 K-向量空间 \(V\) 的维数为 1（\(n^{\circ}\) 2, 命题 4 的推论 1）；把 \(\Omega\) 与 \(V\) 的一个子 A-模等同起来。

命题 8. A-模 V/Ω 是 Matlis 模。

考虑正合列

\[
0 \to \Omega \to \mathrm{V} \to \mathrm{V} / \Omega \to 0;
\]

A-模 V 是内射的（A, X, 第 18 页，例 1），且有 \(\mathrm{di}_{\mathrm{A}}(\Omega) = 1\)（\(n^{\circ} 1\)，命题 1）。由此一方面推出 \(V/\Omega\) 是内射的（\(\S 3\)，\(n^{\circ} 1\)，命题 1），另一方面推出：对 A 的每个极大理想 \(\mathfrak{m}\)，\(\Lambda/\mathfrak{m}\)-向量空间 \(\mathrm{Hom}_{\mathrm{A}}(\mathrm{A}/\mathfrak{m}, \mathrm{V}/\Omega)\) 同构于 \(\mathrm{Ext}_{\mathrm{A}}^{1}(\mathrm{A}/\mathfrak{m}, \Omega)\)，故维数为 1。由于 \(V/\Omega\) 是扭模，它的伴随素理想都是极大的；这就证明了命题（\(\S 8\)，\(n^{\circ} 4\)）。

设 M 为 A-模；依照前引文，用 D(M) 表示 A-模 \(\mathrm{Hom}_{\mathrm{A}}(\mathrm{M},\mathrm{V}/\Omega)\)。可以把第 5 节的构造应用于此，取 I 为复形

\[
\dots 0 \longrightarrow \mathrm{V} \xrightarrow {p} \mathrm{V} / \Omega \longrightarrow 0 \dots
\]

其中 V 置于次数 0，而 p 表示典范满射。复形 D(M) 为

\[
\dots 0 \longrightarrow \operatorname{Hom} _ {\mathrm{A}} (\mathrm{M}, \mathrm{V}) \xrightarrow {p _ {\mathrm{M}}} \mathrm{D} (\mathrm{M}) \longrightarrow 0 \dots ,
\]

其中 \(p_{\mathbf{M}} = \mathrm{Hom}(\mathbf{1}_{\mathbf{M}}, p)\)。我们有分次 A-模的一个典范同构 \(\varphi(\mathbf{M}, \mathbf{I}) : \mathrm{H}(\mathbf{D}(\mathbf{M})) \longrightarrow \mathrm{Ext}_{\mathbf{A}}(\mathbf{M}, \Omega)\)（A, X, 第 100 页，定理 1）。

当 M 是扭模时，模 \(\mathbf{D}(\mathbf{M})^{0} = \operatorname{Hom}_{\mathbf{A}}(\mathbf{M}, \mathbf{V})\) 为零，且 \(\varphi(\mathbf{M}, \mathbf{I})\) 是从 D(M) 到 \(\operatorname{Ext}_{\Lambda}^{1}(\mathbf{M}, \Omega)\) 的同构；态射 \(\alpha_{M}: M \to \mathbf{D}(\mathbf{D}(\mathbf{M}))\) 不是别的，正是 A-模的典范同态

\[
\alpha_ {\mathrm{M}}: \mathrm{M} \longrightarrow \mathrm{D} (\mathrm{D(M)})
\]

它定义于 § 8, 第 3 节，且当 M 是有限型的（即有限长度的）时是同构。在此情形我们重新得到前引文中的状况。

回到一般情形，并设 A-模 M 是有限型的。于是 D(M) 是扭模，故 A-模 \(\mathbf{D}(\mathbf{D}(\mathbf{M}))^{-1} = \mathrm{Hom}_{\mathrm{A}}(\mathrm{D}(\mathrm{M}), \mathrm{V})\) 为零。另一方面，A-模 \(\mathrm{Hom}_{\mathrm{A}}(\mathbf{D}(\mathbf{M})^{0}, \mathrm{V}) = \mathrm{Hom}_{\mathrm{A}}(\mathrm{Hom}_{\mathrm{A}}(\mathrm{M}, \mathrm{V}), \mathrm{V})\) 自然等同于 \(K \otimes_{A} M\)，于是同态

\[
\operatorname{Hom} (1, p): \operatorname{Hom} _ {\mathrm{A}} (\operatorname{Hom} _ {\mathrm{A}} (\mathrm{M}, \mathrm{V}), \mathrm{V}) \longrightarrow \mathrm{D} (\operatorname{Hom} _ {\mathrm{A}} (\mathrm{M}, \mathrm{V}))
\]

等同于映射

\[
j: \mathrm{K} \otimes_ {\mathrm{A}} \mathrm{M} \longrightarrow \mathrm{D} (\operatorname{Hom} _ {\mathrm{A}} (\mathrm{M}, \mathrm{V}))
\]

它满足 \(j(\lambda \otimes m)(f) = p(\lambda f(m))\)（对 \(\lambda \in \mathbf{K}\)，\(m \in \mathbf{M}\)，\(f \in \mathrm{Hom}_{\mathrm{A}}(\mathrm{M},\mathrm{V})\)）。于是第 5 节的定理 1 就转化为序列的正合性

\[
0 \longrightarrow \mathrm{M} \xrightarrow {(i , \alpha_ {\mathrm{M}})} (\mathrm{K} \otimes_ {\mathrm{A}} \mathrm{M}) \oplus \mathrm{D} (\mathrm{D} (\mathrm{M})) \xrightarrow {(j , - \mathrm{D} (p _ {\mathrm{M}}))} \mathrm{D} (\operatorname{Hom} _ {\Lambda} (\mathrm{M}, \mathrm{V})) \longrightarrow 0,
\]

其中 \(i\) 表示从 M 到 \(\mathrm{K} \otimes_{\mathrm{A}} \mathrm{M}\) 的典范映射。\(i\) 的核等同于 M 的扭子模 \(\mathrm{T(M)}\)，它的余核等同于 \((\mathrm{K/A}) \otimes_{\mathrm{A}} \mathrm{M}\)。考虑行正合的交换图表

\[
\begin{array}{c c c c c c c} 0 \to \mathrm{T(M)} & \longrightarrow & \mathrm{M} & \stackrel {{\imath}} {{\longrightarrow}} & \mathrm {K\otimes_ {A} M} & \longrightarrow & (\mathrm{K/A}) \otimes_ {\mathrm{A}} \mathrm{M} \longrightarrow 0 \\ & & ^ {\alpha_ {\mathrm{M}}} \Bigg \downarrow & & \Bigg \downarrow_ {j} \\ 0 \to \mathrm{D(Ext} _ {\mathrm{A}} ^ {1} (\mathrm{M}, \Omega)) & \longrightarrow & \mathrm{D(D(M))} & \stackrel {{\mathrm{D} (p _ {\mathrm{M}})}} {{\longrightarrow}} & \mathrm{D(Hom} _ {\mathrm{A}} (\mathrm{M}, \mathrm{V})) & \longrightarrow & \mathrm{D(Hom} _ {\mathrm{A}} (\mathrm{M}, \Omega)) \longrightarrow 0 \end{array}
\]

其中第二行是由正合列经 Matlis 对偶得到的

\[
0 \rightarrow \operatorname{Hom} _ {\mathrm{A}} (\mathrm{M}, \Omega) \longrightarrow \operatorname{Hom} _ {\mathrm{A}} (\mathrm{M}, \mathrm{V}) \xrightarrow {p _ {\mathrm{M}}} \operatorname{Hom} _ {\mathrm{A}} (\mathrm{M}, \mathrm{V} / \Omega) \longrightarrow \operatorname{Ext} _ {\Lambda} ^ {1} (\mathrm{M}, \Omega) \rightarrow 0.
\]

于是定理 1 的含义是：A-模的同态

\[
\gamma^ {0} (\mathrm{M}): \mathrm{T} (\mathrm{M}) \longrightarrow \mathrm{D} \left(\operatorname{Ext} _ {\mathrm{A}} ^ {1} (\mathrm{M}, \boldsymbol {\Omega})\right) e t \quad \gamma^ {1} (\mathrm{M}): (\mathrm{K} / \mathrm{A}) \otimes_ {\Lambda} \mathrm{M} \longrightarrow \mathrm{D} \left(\operatorname{Hom} _ {\mathrm{A}} (\mathrm{M}, \boldsymbol {\Omega})\right)
\]

——它们分别由 \(\alpha_{\mathrm{M}}\) 与 \(j\) 导出——是双射的。由于 A-模 T(M) 是有限长度的，A-模 \(\operatorname{Ext}_{\mathrm{A}}^{1}(\mathrm{M},\Omega)\) 是有限长度的，且等同于 Matlis 对偶 D(T(M))，并且在群 Z0(A) 中有 \([\operatorname{Ext}_{\mathrm{A}}^{1}(\mathrm{M},\Omega)] = [\operatorname{T}(\mathrm{M})]\) 以及 \(\operatorname{long}_{\mathrm{A}}(\operatorname{Ext}_{\mathrm{A}}^{1}(\mathrm{M},\Omega)) = \operatorname{long}_{\mathrm{A}}(\operatorname{T}(\mathrm{M}))\)（\(\S 8\)，第 4 节，命题 4）。另一方面，取 M = A，便得到典范同构 \(\gamma^{1}(\mathrm{A}):K / A\to D(\Omega)\)。

设 B 为 K 的一个包含 A 的子环，且在 \(\Lambda\) 上是有限的。对 A 的每个极大理想 \(\mathfrak{m}\)，有 \(\mathrm{prof}_{\mathrm{A}_{\mathfrak{m}}}(B_{\mathfrak{m}}) = \dim_{A_{\mathfrak{m}}}(B_{\mathfrak{m}}) = 1\)（\(\S 1\)，\(n^{\circ} 1\)，注 2 及 VIII, \(\S 2\)，\(n^{\circ} 3\)，定理 1），故 B 是 Macaulay A-模。于是，B-模 \(\Omega_{\mathrm{B}} = \mathrm{Hom}_{\mathrm{A}}(B, \Omega)\) 是对偶化的（\(n^{\circ} 3\)，注）。从 \(\Omega_{\mathrm{B}} = \mathrm{Hom}_{\mathrm{A}}(B, \Omega)\) 到 \(\Omega = \mathrm{Hom}_{\mathrm{A}}(A, \Omega)\) 的典范映射是单射的；它的像由 \(\omega\) 的这样的元素 \(\Omega\) 组成：V 的子 B-模 B\(\omega\) 含于 \(\Omega\)。于是 \(\Omega_{\mathrm{B}}\) 等同于 \(\Omega\) 的最大子 B-模。A-模 B/\(\Lambda\) 是有限长度的；正合列

\[
0 \to \Omega_ {\mathrm{B}} \to \Omega \to \operatorname{Ext} _ {\mathrm{A}} ^ {1} (\mathrm{B/A}, \Omega) \to 0
\]

使我们可以把 \(\Omega/\Omega_{\mathrm{B}}\) 与 \(\operatorname{Ext}_{\mathrm{A}}^{1}(\mathrm{B}/\mathrm{A},\Omega)\) 等同，从而由前所述与 \(\mathrm{D}(\mathrm{B}/\mathrm{A})\) 等同。特别地，在 \([\mathrm{B}/\mathrm{A}] = [\Omega/\Omega_{\mathrm{B}}]\) 中有 \(Z_{0}(\mathrm{A})\) 以及 \(\operatorname{Ann}_{\mathrm{A}}(\mathrm{B}/\mathrm{A}) = \operatorname{Ann}_{\mathrm{A}}(\Omega/\Omega_{\mathrm{B}})\)（\(\S 8\)，\(\mathfrak{n}^{\circ}4\)，命题 4）。

理想 \(\mathfrak{c} = \mathrm{Ann}_{\mathrm{A}}(\mathrm{B / A})\) 是传递子 A : B，即（VII, § 1, 第 1 节）K 的元素 \(x\)（满足 \(xB \subset \Lambda\)）的集合。它是

A 的（非零）理想，也是 B 的理想；事实上它是 B 的含于 A 的最大理想。由于 \(\Omega_{\mathrm{B}} : \Omega \subset \Omega : \Omega = \mathrm{A} (\mathrm{n}^{\circ} 4, \text{prop. } 7, \mathrm{b}))\)，有 Ann \(_{\mathrm{A}}(\Omega/\Omega_{\mathrm{B}}) = \Omega_{\mathrm{B}} : \Omega\)，最终得到

\[
\mathfrak {c} = \operatorname{Ann} _ {\mathrm{A}} (\mathrm{B/A}) = \operatorname{Ann} _ {\mathrm{A}} (\Omega / \Omega_ {\mathrm{B}}) = \Omega_ {\mathrm{B}}: \Omega .
\]

由于 \(\Omega_{\mathrm{B}}\) 是 B-模，关系 \(x\Omega \subset \Omega_{\mathrm{B}}\) 等价于 \(xB\Omega \subset \Omega_{\mathrm{B}}\)，故又有 \(\mathfrak{c} = \Omega_{\mathrm{B}}: \mathrm{B}\Omega\)。

现在我们把前面的讨论特殊化到 B 是 A 的整闭包的情形；当环 A 是 Japanese 环（IX, § 4, 第 1 节，定义 1）时，B 是有限生成 A-模这一假设得到满足，而当 A 是局部且完备的环（前引文，第 2 节，定理 2），或基本上是在域上有限型的（前引文，第 1 节，注 2 及例）时就属于这种情形。此时环 B 是 Dedekind 环（VII, § 2, 第 2 节，定理 1），且无扭的 B-模 \(\Omega_{\mathrm{B}}\)、B \(\Omega\) 与 c 都是秩为 1 的投射模（VII, § 4, 第 10 节，命题 22）。于是关系 \(c = \Omega_{\mathrm{B}} : \mathrm{B}\Omega\) 的含义是：由 K 在 V 上的作用导出的线性映射 \(c \otimes_{\mathrm{B}} \mathrm{B}\Omega \to \Omega_{\mathrm{B}}\) 是同构（II, § 5, 第 6 节，命题 11）。特别地，有 \(\Omega_{\mathrm{B}} = c(\mathrm{B}\Omega) = c\Omega\)。

命题 9. 设 B 为 A 的整闭包，并令 \(\mathfrak{c} = \mathrm{A}:\mathrm{B}\)。假设 B 是有限生成的 A-模。在 \([\mathrm{B} / \mathrm{c}] \leqslant 2[\mathrm{B} / \mathrm{A}]\) 中有不等式 \(\mathrm{Z}_0(\mathrm{A})\)。为使等号成立，必须且只须 A 是 Gorenstein 环。

我们有 \([B/c] = [B/A] + [A/c]\)，故所考虑的不等式等价于 \([A/c] \leqslant [B/A]\)。

\begin{enumerate}
\def\labelenumi{\Alph{enumi})}
\item
  对 A 的每个极大理想 m，\(A_{m}\) 的整闭包是 \(B_{m}\)（V, § 1, 第 5 节，推论 1），且有 \(c_{m} = A_{m} : B_{m}\)。此外，按定义有 \([B/c] = \sum_{m} \text{long}_{A_{m}}(B_{m}/c_{m})[m]\) 与 \([B/A] = \sum_{m} \text{long}_{A_{m}}(B_{m}/A_{m})[m]\)。这就化归为环 A 是局部环的情形来证明命题，以下即作此假设。在此情形，有序群 \(Z_{0}(A)\) 典范地等同于 Z，于是有限长度模的类就是它的长度。环 B 是半局部的，且 B-模 BΩ 是秩为 1 的自由模（II, § 5, 第 3 节，命题 5）。
\item
  若 A 是 Gorenstein 环，则 A-模 \(\Omega\) 是秩为 1 的自由模（第 4 节，命题 7 的推论 1）；选取 \(\omega\) 的一个生成元 \(\Omega\)。有 \(\Omega = A\omega\) 与 \(\Omega_{\mathrm{B}} = c\Omega = c\omega\)，从而 \(\text{long}(A/c) = \text{long}(\Omega/\Omega_{\mathrm{B}}) = \text{long}(B/A)\)。
\item
  设剩余域 \(\kappa_{\mathrm{A}}\) 是无限的。对 B 的每个极大理想 \(\mathfrak{n}\)，用 L(n) 表示 B/n-向量空间（一维的）BΩ/nBΩ，并用 pr\(_{\mathfrak{n}}\) 表示从 \(\bigoplus_{\mathfrak{n}}\) L(n) 到 L(n) 上的典范投影。设 \(\varphi : \Omega \longrightarrow \bigoplus_{\mathfrak{n}}\) L(n) 是 \(\Omega\) 经典范同态 BΩ \(\longrightarrow \bigoplus_{\mathfrak{n}}\) L(n) 而得到的映射。\(\varphi\) 的像是 \(\kappa_{\mathrm{A}}\)-向量空间 \(\bigoplus_{\mathfrak{n}}\) L(n) 的一个子空间；它不含于 Ker pr \(_{\mathfrak{n}}\)，因为否则将有 \(\Omega \subset \mathfrak{n}\mathrm{B}\Omega\)，从而 BΩ \(\subset \mathfrak{n}\mathrm{B}\Omega\)，这是矛盾的。于是 \(\varphi\) 的像不含于诸 Ker pr \(_{\mathfrak{n}}\) 的并（A, V, 第 40 页，引理 1）；因此存在一个元素 \(\omega\)（属于 \(\Omega\)），它在 BΩ/nBΩ 中的像对每个 n 都非零，这蕴含 \(\omega\) 生成 B-模 BΩ（II, § 3, 第 3 节，命题 11）。
\end{enumerate}

设 \(a \in \mathbf{A}\)；若 \(a\omega\) 属于 \(\Omega_{\mathrm{B}}\)，则有 \(a\mathrm{B}\omega \subset \Omega_{\mathrm{B}}\)，从而 \(a\Omega \subset \Omega_{\mathrm{B}}\)，这蕴含 \(a \in \mathfrak{c}\)。于是映射 \(a \mapsto a\omega\) 诱导出从 \(\mathrm{A}/\mathfrak{c}\) 到 \(\Omega/\Omega_{\mathrm{B}}\) 的一个单射；从而有 \(\mathrm{long}(\mathrm{A}/\mathfrak{c}) \leqslant \mathrm{long}(\Omega/\Omega_{\mathrm{B}}) = \mathrm{long}(\mathrm{B}/\mathrm{A})\)。

若 long(A/c) = long(B/A)，则 Aω + Ω＿B = Ω。可设理想 c 含于 m＿A（否则 A 等于 B，从而是 Gorenstein 环）。由于 Ω＿B = cΩ 含于 m＿AΩ，由 Nakayama 引理，ω 生成 Ω。于是 A-模 Ω 是循环模，从而是秩为 1 的自由模，这意味着 A 是 Gorenstein 环（第 4 节，命题 7 的推论 1）。

\begin{enumerate}
\def\labelenumi{\Alph{enumi})}
\setcounter{enumi}{3}
\tightlist
\item
  现在处理一般情形。用 A′ 表示环 A{]}X{[}，即（IX, App., 第 2 节）多项式环 Λ{[}X{]} 在素理想 m\_AΛ{[}X{]} 处的局部环；它是一维的平坦整 A-代数，其剩余域 κ\_A′ 等同于 κ\_A(X)，而分式域等同于 K(X)（前引文）。由第 3 节命题 5 的推论，A′-模 A′ ⊗\_A Ω 是对偶化的。令 B′ = A′ ⊗\_A B；它是 Λ′ 在 K(X) 中的整闭包（V, § 1, 第 3 节，命题 13 及第 5 节，命题 16）。传递子 c′ = A′ : B′ 等于 cA′（I, § 2, 第 10 节，公式 (11)）。对每个有限长度的 A-模 M，有 long\_A′(A′ ⊗\_A M) = long\_A(M)：事实上，由于 A-代数 A′ 是平坦的，只需对 M 是单模（即同构于 κ\_A）的情形证明此关系；但在这种情形，A′ ⊗\_A κ\_A 等同于 κ\_A′，断言得证。于是我们有
\end{enumerate}

\[
\operatorname{long} _ {\mathrm{A}} (\mathrm{B} / \mathfrak {c}) = \operatorname{long} _ {\mathrm{A} ^ {\prime}} \left(\mathrm{B} ^ {\prime} / \mathfrak {c} ^ {\prime}\right) \quad \text { and } \quad \operatorname{long} _ {\mathrm{A}} (\mathrm{B} / \mathrm{A}) = \operatorname{long} _ {\mathrm{A} ^ {\prime}} \left(\mathrm{B} ^ {\prime} / \mathrm{A} ^ {\prime}\right).
\]

环 A′ 满足命题的假设，且它的剩余域是无限的。由证明的 C) 部分，有 \(\text{long}_{\text{A}'}(\text{B}'/\text{c}') \leqslant 2 \text{long}_{\text{A}'}(\text{B}'/\text{A}')\)，且等号成立蕴含 A′ 是 Gorenstein 环；而后者又蕴含 A 是 Gorenstein 环（§ 3, 第 8 节，命题 12 的推论 1）。

\section{局部上同调，Grothendieck 对偶}
